\documentclass[fs11,seriesFre]{korfarm-base}
\usepackage{korfarm-frege}

% ============================================================
% passage-note.tex (v11)
% 지문 박스 + 첫 페이지 우측 메모란 (동적 높이)
% 정책: PAGE_BREAK_POLICY.md §3, §4 참조
%
% 변경 (v11):
%  · 지문 박스 폭 확대: enlarge right by=-3.7cm  (메모 3.4cm + gap 3mm)
%  · 메모란 동적 높이: frame.north east ~ frame.south east 사이 자동 채움
%  · 메모 줄선: clip 영역으로 박스 내부만 채움 (박스 높이 모르고도 작동)
%  · 문단 들여쓰기 활성화: tcolorbox 내부 \parindent=1em, \noindent 제거
%
% 사용:
%   \kfPassageNote{제목}{지문 본문}      → 메모란 포함
%   \kfPassageNoteNT{지문 본문}          → 메모란 포함, 제목 없음
%   \kfPassageFull{제목}{본문}           → 메모란 없음 (활동 내 보기)
% ============================================================

\makeatletter

% 메모란 규격 (cm 고정)
\newlength{\kf@memoW}
\setlength{\kf@memoW}{3.4cm}
\newlength{\kf@memoGap}
\setlength{\kf@memoGap}{3mm}

% --- 메인: 제목 있는 버전 ---
\newcommand{\kf@passagenote@with}[2]{%
  \par\ifkfPassageNewPage\ifdim\pagetotal>0.4\textheight\clearpage\fi\fi\smallskip\needspace{6\baselineskip}%
  \begin{tcolorbox}[
    enhanced, breakable,
    width=\dimexpr\linewidth-4cm\relax,
    colback=kfPassageBg, colframe=kfPrimary!70,
    boxrule=0.6pt, arc=3pt,
    left=\kfBoxPad, right=\kfBoxPad, top=\dimexpr\kfBoxPad+8pt\relax, bottom=\kfBoxPad,
    title={\small\bfseries\color{kfPrimary} #1},
    coltitle=kfPrimary,
    colbacktitle=kfPrimaryLight!60,
    attach boxed title to top left={yshift=-2.4mm, xshift=8pt},
    boxed title style={colframe=kfPrimary!50, boxrule=0.5pt, arc=2pt},
    parbox=false,
    overlay unbroken and first={%
      \begin{scope}
        \draw[kfRule, line width=0.4pt, fill=kfNoteBg, rounded corners=2pt]
          ([xshift=\kf@memoGap]frame.north east) rectangle
          ([xshift=\kf@memoGap+\kf@memoW]frame.south east);
        \node[anchor=north east, font=\footnotesize\bfseries\color{kfMute}, inner sep=3pt]
          at ([xshift=\kf@memoGap+\kf@memoW, yshift=-2pt]frame.north east) {메모};
        \node[anchor=north east, fill=kfPrimary!15, text=kfPrimary,
              font=\scriptsize\bfseries, rounded corners=2pt,
              inner xsep=4pt, inner ysep=2pt]
          at ([xshift=\kf@memoGap+\kf@memoW-3pt, yshift=-19pt]frame.north east) {\kf@memocat};
        \begin{scope}
          \path[clip]
            ([xshift=\kf@memoGap]frame.north east) rectangle
            ([xshift=\kf@memoGap+\kf@memoW]frame.south east);
          \foreach \i in {1,...,40}{%
            \draw[kfRule, line width=0.25pt]
              ([xshift=\kf@memoGap+2mm, yshift=-7mm - \i*5mm]frame.north east) --
              ([xshift=\kf@memoGap+\kf@memoW-2mm, yshift=-7mm - \i*5mm]frame.north east);%
          }
        \end{scope}
      \end{scope}
    }
  ]
    \setlength{\parindent}{1em}%
    \hspace*{1em}\ignorespaces#2%
  \end{tcolorbox}\par\smallskip
}

% --- 제목 없는 버전 ---
\newcommand{\kf@passagenote@notitle}[1]{%
  \par\ifkfPassageNewPage\ifdim\pagetotal>0.4\textheight\clearpage\fi\fi\smallskip\needspace{6\baselineskip}%
  \begin{tcolorbox}[
    enhanced, breakable,
    width=\dimexpr\linewidth-4cm\relax,
    colback=kfPassageBg, colframe=kfPrimary!70,
    boxrule=0.6pt, arc=3pt,
    left=\kfBoxPad, right=\kfBoxPad, top=\dimexpr\kfBoxPad+8pt\relax, bottom=\kfBoxPad,
    parbox=false,
    overlay unbroken and first={%
      \begin{scope}
        \draw[kfRule, line width=0.4pt, fill=kfNoteBg, rounded corners=2pt]
          ([xshift=\kf@memoGap]frame.north east) rectangle
          ([xshift=\kf@memoGap+\kf@memoW]frame.south east);
        \node[anchor=north east, font=\footnotesize\bfseries\color{kfMute}, inner sep=3pt]
          at ([xshift=\kf@memoGap+\kf@memoW, yshift=-2pt]frame.north east) {메모};
        \node[anchor=north east, fill=kfPrimary!15, text=kfPrimary,
              font=\scriptsize\bfseries, rounded corners=2pt,
              inner xsep=4pt, inner ysep=2pt]
          at ([xshift=\kf@memoGap+\kf@memoW-3pt, yshift=-19pt]frame.north east) {\kf@memocat};
        \begin{scope}
          \path[clip]
            ([xshift=\kf@memoGap]frame.north east) rectangle
            ([xshift=\kf@memoGap+\kf@memoW]frame.south east);
          \foreach \i in {1,...,40}{%
            \draw[kfRule, line width=0.25pt]
              ([xshift=\kf@memoGap+2mm, yshift=-7mm - \i*5mm]frame.north east) --
              ([xshift=\kf@memoGap+\kf@memoW-2mm, yshift=-7mm - \i*5mm]frame.north east);%
          }
        \end{scope}
      \end{scope}
    }
  ]
    \setlength{\parindent}{1em}%
    \hspace*{1em}\ignorespaces#1%
  \end{tcolorbox}\par\smallskip
}

\newcommand{\kfPassageNote}[2]{%
  \def\kf@tmptitle{#1}%
  \ifx\kf@tmptitle\@empty
    \kf@passagenote@notitle{#2}%
  \else
    \kf@passagenote@with{#1}{#2}%
  \fi
}
\newcommand{\kfPassageNoteNT}[1]{\kf@passagenote@notitle{#1}}
\makeatother

% ============================================================
% 풀폭 지문 박스 (메모란 없음) — 활동 내 보기 박스
% PAGE_BREAK_POLICY.md §3
% ============================================================
\makeatletter
\newcommand{\kf@passagefull@with}[2]{%
  \par\ifkfPassageNewPage\ifdim\pagetotal>0.4\textheight\clearpage\fi\fi\smallskip\needspace{6\baselineskip}%
  \begin{tcolorbox}[
    enhanced, breakable,
    colback=kfPassageBg, colframe=kfPrimary!70,
    boxrule=0.6pt, arc=3pt,
    left=\kfBoxPad, right=\kfBoxPad, top=\dimexpr\kfBoxPad+8pt\relax, bottom=\kfBoxPad,
    title={\small\bfseries\color{kfPrimary} #1},
    coltitle=kfPrimary,
    colbacktitle=kfPrimaryLight!60,
    attach boxed title to top left={yshift=-2.4mm, xshift=8pt},
    boxed title style={colframe=kfPrimary!50, boxrule=0.5pt, arc=2pt},
    parbox=false
  ]
    \setlength{\parindent}{1em}%
    \hspace*{1em}\ignorespaces#2%
  \end{tcolorbox}\par\smallskip
}
\newcommand{\kf@passagefull@notitle}[1]{%
  \par\ifkfPassageNewPage\ifdim\pagetotal>0.4\textheight\clearpage\fi\fi\smallskip\needspace{6\baselineskip}%
  \begin{tcolorbox}[
    enhanced, breakable,
    colback=kfPassageBg, colframe=kfPrimary!70,
    boxrule=0.6pt, arc=3pt,
    left=\kfBoxPad, right=\kfBoxPad, top=\dimexpr\kfBoxPad+8pt\relax, bottom=\kfBoxPad,
    parbox=false
  ]
    \setlength{\parindent}{1em}%
    \hspace*{1em}\ignorespaces#1%
  \end{tcolorbox}\par\smallskip
}
\newcommand{\kfPassageFull}[2]{%
  \def\kf@tmptitle{#1}%
  \ifx\kf@tmptitle\@empty
    \kf@passagefull@notitle{#2}%
  \else
    \kf@passagefull@with{#1}{#2}%
  \fi
}
\newcommand{\kfPassageFullNT}[1]{\kf@passagefull@notitle{#1}}
\makeatother

% ============================================================
% 작은 문장 박스 (문장 독해용) — PAGE_BREAK_POLICY.md §2.5
% ============================================================
\newcommand{\kfSentenceBox}[2]{%
  \par\smallskip\needspace{3\baselineskip}%
  \noindent\begin{tcolorbox}[
    enhanced, breakable=false,
    colback=kfPassageBg, colframe=kfMute,
    boxrule=0.4pt, arc=2pt,
    left=8pt, right=6pt, top=4pt, bottom=4pt,
    title={\footnotesize\bfseries\color{kfPrimary} #1},
    coltitle=kfPrimary, colbacktitle=kfPaper,
    attach boxed title to top left={yshift=-1.6mm, xshift=4pt},
    boxed title style={colframe=kfMute, boxrule=0.3pt, arc=1pt}
  ]%
    \setstretch{1.3}#2%
  \end{tcolorbox}\par\smallskip%
}

% ============================================================
% 지문 본문 아래 안내/정리 박스 (v16 신규)
% 사용: \kfPassageHint{한 줄 정리}{본문 안내 텍스트}
% ============================================================
\newcommand{\kfPassageHint}[2]{%
  \par\smallskip\needspace{4\baselineskip}%
  \begin{tcolorbox}[
    enhanced, breakable=false,
    colback=kfPrimary!5, colframe=kfPrimary!40,
    boxrule=0.4pt, arc=2pt,
    left=8pt, right=6pt, top=4pt, bottom=4pt,
    title={\footnotesize\bfseries\color{kfPrimary} #1},
    coltitle=kfPrimary, colbacktitle=kfPaper,
    attach boxed title to top left={yshift=-1.5mm, xshift=4pt},
    boxed title style={colframe=kfPrimary!40, boxrule=0.3pt, arc=1pt}
  ]%
    \small\setstretch{1.3}#2%
  \end{tcolorbox}\par\smallskip%
}

% ============================================================
% 환경 버전 (호환성) — 메모란 포함
% ============================================================
\makeatletter
\newenvironment{kfPassageNoteEnv}[1]%
  {\par\needspace{6\baselineskip}%
   \def\kf@psrc{#1}%
   \begin{tcolorbox}[
     enhanced, breakable,
     width=\dimexpr\linewidth-4cm\relax,
     colback=kfPassageBg, colframe=kfPrimary!70,
     boxrule=0.6pt, arc=3pt,
     left=\kfBoxPad, right=\kfBoxPad, top=\dimexpr\kfBoxPad+8pt\relax, bottom=\kfBoxPad,
     title={\small\bfseries\color{kfPrimary} \kf@psrc},
     coltitle=kfPrimary, colbacktitle=kfPrimaryLight!60,
     attach boxed title to top left={yshift=-2.4mm, xshift=8pt},
     boxed title style={colframe=kfPrimary!50, boxrule=0.5pt, arc=2pt},
     parbox=false,
     overlay unbroken and first={%
       \begin{scope}
         \draw[kfRule, line width=0.4pt, fill=kfNoteBg, rounded corners=2pt]
           ([xshift=\kf@memoGap]frame.north east) rectangle
           ([xshift=\kf@memoGap+\kf@memoW]frame.south east);
         \node[anchor=north east, font=\footnotesize\bfseries\color{kfMute}, inner sep=3pt]
           at ([xshift=\kf@memoGap+\kf@memoW, yshift=-2pt]frame.north east) {메모};
         \begin{scope}
           \path[clip]
             ([xshift=\kf@memoGap]frame.north east) rectangle
             ([xshift=\kf@memoGap+\kf@memoW]frame.south east);
           \foreach \i in {1,...,40}{%
             \draw[kfRule, line width=0.25pt]
               ([xshift=\kf@memoGap+2mm, yshift=-7mm - \i*5mm]frame.north east) --
               ([xshift=\kf@memoGap+\kf@memoW-2mm, yshift=-7mm - \i*5mm]frame.north east);%
           }
         \end{scope}
       \end{scope}
     }
   ]%
   \setlength{\parindent}{1em}%
  }%
  {\end{tcolorbox}\par\smallskip}
\makeatother

% ============================================================
% condition-box.tex
% <조건> 박스 컴포넌트 (tcolorbox)
% 사용:
%   \begin{kfCondition}
%     \item 첫째 조건
%     \item 둘째 조건
%   \end{kfCondition}
% ============================================================

\newenvironment{kfCondition}%
  {\par\smallskip\needspace{4\baselineskip}%
   \begin{tcolorbox}[
     enhanced, breakable=false,
     colback=kfPrimaryLight!40, colframe=kfPrimary,
     boxrule=0.5pt, arc=2pt,
     left=\kfBoxPad, right=\kfBoxPad, top=\kfBoxPad, bottom=\kfBoxPad,
     title={\small\bfseries\color{kfPrimary} \,조건\,},
     coltitle=kfPrimary, colbacktitle=kfPaper,
     attach boxed title to top left={yshift=-2mm, xshift=6pt},
     boxed title style={colframe=kfPrimary, boxrule=0.4pt, arc=1pt}
   ]%
   \begin{enumerate}[leftmargin=16pt, itemsep=1pt, topsep=1pt,
                    label=\textbf{\arabic*.}]}%
  {\end{enumerate}\end{tcolorbox}\par\smallskip}

% ============================================================
% evidence-box.tex
% <보기> 박스 컴포넌트
% 사용:
%   \begin{kfEvidence}
%     보기 본문 ...
%   \end{kfEvidence}
% ============================================================

\newenvironment{kfEvidence}%
  {\par\smallskip\needspace{4\baselineskip}%
   \begin{tcolorbox}[
     enhanced, breakable=false,
     colback=kfPaper, colframe=kfMute,
     boxrule=0.5pt, arc=2pt,
     left=\kfBoxPad, right=\kfBoxPad, top=\kfBoxPad, bottom=\kfBoxPad,
     title={\small\bfseries\color{kfInk} \,보기\,},
     coltitle=kfInk, colbacktitle=kfPrimaryLight!50,
     attach boxed title to top left={yshift=-2mm, xshift=6pt},
     boxed title style={colframe=kfMute, boxrule=0.4pt, arc=1pt}
   ]%
   \leavevmode\mbox{}\ignorespaces%
  }%
  {\end{tcolorbox}\par\smallskip}

% ============================================================
% choice-list.tex
% 5선지 (① ② ③ ④ ⑤) 자동 들여쓰기 컴포넌트
% 사용:
%   \begin{kfChoices}
%     \kfChoice 첫째 선택지
%     \kfChoice 둘째 선택지
%     \kfChoice 셋째 선택지
%     \kfChoice 넷째 선택지
%     \kfChoice 다섯째 선택지
%   \end{kfChoices}
% ============================================================

\newcounter{kfChoiceCnt}
\newcommand{\kfChoiceMark}[1]{%
  \ifcase#1\relax%
  \or ①\or ②\or ③\or ④\or ⑤\or ⑥\or ⑦\or ⑧\or ⑨%
  \fi
}

% 환경: 자동 번호
\newenvironment{kfChoices}%
  {\setcounter{kfChoiceCnt}{0}%
   \par\smallskip%
   \begin{list}{}{%
     \setlength{\leftmargin}{20pt}%
     \setlength{\itemindent}{0pt}%
     \setlength{\labelwidth}{18pt}%
     \setlength{\labelsep}{6pt}%
     \setlength{\itemsep}{2pt}%
     \setlength{\topsep}{1pt}%
     \setlength{\parsep}{0pt}%
     \renewcommand{\makelabel}[1]{##1\hss}%
   }}%
  {\end{list}\par\smallskip}

\newcommand{\kfChoice}{%
  \stepcounter{kfChoiceCnt}%
  \item[\textbf{\kfChoiceMark{\value{kfChoiceCnt}}}]\ignorespaces%
}

% --- 한 줄 5선지 (짧은 선택지용) ---
\newcommand{\kfChoicesInline}[5]{%
  \par\smallskip%
  \noindent\textbf{①}~#1\quad\textbf{②}~#2\quad\textbf{③}~#3\quad\textbf{④}~#4\quad\textbf{⑤}~#5%
  \par\smallskip%
}

% ============================================================
% answer-note.tex
% 답안 작성란 (노트 스타일, 줄친 영역)
% 사용:
%   \kfAnswerLines{5}        % 5줄 답안란
%   \kfAnswerNote{서술형 답안란}{8} % 라벨+8줄
% ============================================================

% 단일 줄 (필요 시 인라인)
\newcommand{\kfAnswerLine}{%
  \par\noindent\textcolor{kfRule}{\rule{\linewidth}{0.4pt}}\par
}

% N줄 답안란 (간단)
\newcommand{\kfAnswerLines}[1]{%
  \par\smallskip%
  \begin{minipage}{\linewidth}%
  \foreach \i in {1,...,#1}{%
    \textcolor{kfRule}{\rule{\linewidth}{0.4pt}}\\[6pt]%
  }%
  \end{minipage}%
  \par\smallskip%
}

% 라벨이 있는 노트형 답안란
\newcommand{\kfAnswerNote}[2]{%
  \par\smallskip\needspace{#2\baselineskip}%
  \begin{tcolorbox}[
    enhanced, breakable=false,
    colback=kfPaper, colframe=kfRule,
    boxrule=0.4pt, arc=2pt,
    left=\kfBoxPad, right=\kfBoxPad, top=\kfBoxPad, bottom=\kfBoxPad,
    title={\footnotesize\bfseries\color{kfMute} #1},
    coltitle=kfMute, colbacktitle=kfPaper,
    attach boxed title to top left={yshift=-1.5mm, xshift=6pt},
    boxed title style={colframe=kfRule, boxrule=0.3pt, arc=1pt}
  ]%
  \foreach \i in {1,...,#2}{%
    \textcolor{kfRule}{\rule{\linewidth}{0.3pt}}\\[6pt]%
  }%
  \end{tcolorbox}\par\smallskip%
}

% 짧은 빈칸 (단답형)
\newcommand{\kfBlank}[1][3cm]{\underline{\hspace{#1}}}
% 넓은 빈칸 (서술 한 줄)
\newcommand{\kfBlankWide}[1][6cm]{\underline{\hspace{#1}}}
% 괄호 빈칸
\newcommand{\kfBlankParen}{(\,\hspace{1cm}\,)}
% O/X 빈칸
\newcommand{\kfBlankOX}{(\,\hspace{1.5cm}\,)}

% ============================================================
% question-block.tex
% 문제 블록 (번호 배지 + 발문 + 보기/조건 + 선택지 + 답안란)
% 모든 구성을 하나의 minipage로 묶어 단/페이지 단절 금지
%
% 사용:
%   \begin{kfQBlock}{1}{객관식}
%     \kfQStem{다음 글에 대한 설명으로 가장 적절한 것은?}
%     \begin{kfEvidence} ... \end{kfEvidence}    % 선택
%     \begin{kfCondition} ... \end{kfCondition}  % 선택
%     \begin{kfChoices}
%       \kfChoice ...
%     \end{kfChoices}
%     \kfAnswerLines{2}                            % 선택
%   \end{kfQBlock}
% ============================================================

% 무단절 minipage로 묶고, 발문/보기/조건/선택지/답안란을 자유롭게 배치
% 사용자 피드백 #4: [객관식]/[서술형] 등 텍스트 라벨 제거 — 번호 배지만 표시
% #2(유형 라벨)는 호환성 인자로만 받고 출력하지 않음.
\newenvironment{kfQBlock}[2]%
  {\par\smallskip\needspace{6\baselineskip}%
   \noindent\begin{minipage}{\linewidth}%
   \samepage%
   \par\noindent
   \kfQNum{#1}\hspace{4pt}%
   \begin{minipage}[t]{\dimexpr\linewidth-30pt\relax}%
  }%
  {\end{minipage}\end{minipage}\par\smallskip}

% 문제 발문
\newcommand{\kfQStem}[1]{%
  \par\noindent#1\par\vspace{2pt}%
}

% 짧은 소문항 라벨 (1), (2), (3)
\newcounter{kfSubQ}
\newcommand{\kfSubQReset}{\setcounter{kfSubQ}{0}}
\newcommand{\kfSubQ}{%
  \stepcounter{kfSubQ}%
  \par\noindent\textbf{(\arabic{kfSubQ})}~\ignorespaces%
}

% ============================================================
% activity-table.tex
% 활동: 정리표 / 내용 확인표 (마크다운 표 → tabularx 변환)
% 사용:
%   % 일반 tabular (직접 컬럼 명시)
%   \begin{kfActTable}{|l|X|X|}
%     \kfTHead{항목} & \kfTHead{설명} & \kfTHead{학생 작성} \\\hline
%     ① & ... & \kfTBlank \\\hline
%   \end{kfActTable}
%
%   % adjustbox 자동 축소 (정해진 폭으로 강제)
%   \kfActTableWrap{
%     ... (\begin{tabular}...\end{tabular} 코드) ...
%   }
% ============================================================

% 사용 시 본문 마지막 행은 \\\hline 으로 끝나야 함.
% v17.1: 흰색 mask로 Canva 배경 본문 침범 차단
\newenvironment{kfActTable}[1]%
  {\par\smallskip\needspace{4\baselineskip}%
   \renewcommand{\arraystretch}{1.3}%
   \arrayrulecolor{kfRule}%
   \begin{tcolorbox}[blanker, colback=white, left=2pt, right=2pt, top=2pt, bottom=2pt]%
   \noindent\begin{adjustbox}{max width=\linewidth}%
   \begin{tabular}{#1}%
   \hline}%
  {\end{tabular}%
   \end{adjustbox}%
   \end{tcolorbox}%
   \arrayrulecolor{black}%
   \par\smallskip}

% 헤더 행 헬퍼
\newcommand{\kfTHead}[1]{\textbf{\color{kfPrimary} #1}}
\newcommand{\kfTHeadRow}[1]{\rowcolor{kfPrimaryLight!50}\textbf{\color{kfPrimary} #1}}

% 빈 셀 (학생 작성용)
\newcommand{\kfTBlank}{\rule[-1.6em]{0pt}{3.4em}}
\newcommand{\kfTBlankSm}{\rule[-1.0em]{0pt}{2.4em}}

% adjustbox 강제 축소 래퍼 — Python에서 변환된 raw tabular 블록을 감쌀 때 사용
\newcommand{\kfActTableWrap}[1]{%
  \par\smallskip\needspace{4\baselineskip}%
  \begin{tcolorbox}[blanker, colback=white, left=2pt, right=2pt, top=2pt, bottom=2pt]%
  \noindent\begin{adjustbox}{max width=\linewidth, center}%
  #1%
  \end{adjustbox}%
  \end{tcolorbox}\par\smallskip%
}

% ============================================================
% activity-vocab.tex
% 어휘 학습 표 (3열: 번호 - 뜻 - 빈칸)
% 사용:
%   \begin{kfVocab}
%     \kfVocabItem{1}{눈으로 보고 마음으로 깨달음}
%     \kfVocabItem{2}{사물의 이치를 따져 헤아림}
%   \end{kfVocab}
% ============================================================

% adjustbox로 폭 강제 + 일반 tabular + p{} 열 사용
\newenvironment{kfVocab}%
  {\par\smallskip\needspace{4\baselineskip}%
   \renewcommand{\arraystretch}{1.35}%
   \arrayrulecolor{kfRule}%
   \noindent\begin{adjustbox}{max width=\linewidth, center}%
   \begin{tabular}{|>{\centering\arraybackslash}m{1.0cm}|m{8.5cm}|>{\centering\arraybackslash}m{4.0cm}|}%
   \hline
   \rowcolor{kfPrimaryLight!50}%
   \multicolumn{1}{|c|}{\textbf{\color{kfPrimary}번호}} &
   \multicolumn{1}{c|}{\textbf{\color{kfPrimary}뜻풀이}} &
   \multicolumn{1}{c|}{\textbf{\color{kfPrimary}어휘 (정답 작성)}} \\\hline
   \ignorespaces}%
  {\end{tabular}%
   \end{adjustbox}%
   \arrayrulecolor{black}%
   \par\smallskip}

\newcommand{\kfVocabItem}[2]{%
  \textbf{#1} & #2 & \rule[-1.0em]{0pt}{2.4em}\\\hline%
}

% ============================================================
% activity-blank.tex
% 빈칸 채우기 활동
% 사용:
%   \begin{kfBlankList}
%     \kfBlankItem 첫 번째 문장에서 (\kfBlank)은(는) ...
%     \kfBlankItem 둘째 문장에서 ...
%   \end{kfBlankList}
% ============================================================

\newenvironment{kfBlankList}%
  {\par\smallskip%
   \begin{enumerate}[leftmargin=20pt, itemsep=4pt, topsep=2pt,
                    label=\textbf{\arabic*.}, labelsep=6pt]}%
  {\end{enumerate}\par\smallskip}

\newcommand{\kfBlankItem}{\item}

% 텍스트 안에 박스형 빈칸 (단어 1개 채우기)
\newcommand{\kfBlankBox}[1][2.4cm]{%
  \tikz[baseline=-0.5ex]{%
    \draw[kfRule, line width=0.4pt] (0,0) -- ++(#1,0);%
  }%
}

% ============================================================
% activity-ox.tex (v15)
% O/X 표기 활동 — 그래픽 디자인
% 사용:
%   \begin{kfOXList}
%     \kfOXItem 글쓴이는 자연을 예찬하고 있다.\kfOX
%     \kfOXItem 1연과 4연은 수미상관 구조이다.\kfOX
%   \end{kfOXList}
%
% \kfOX = 우측에 [O] [X] 두 박스 (학생이 하나 동그라미)
% ============================================================

\newenvironment{kfOXList}%
  {\par\smallskip%
   \begin{enumerate}[leftmargin=20pt, itemsep=5pt, topsep=2pt,
                    label=\textbf{\arabic*.}, labelsep=6pt]}%
  {\end{enumerate}\par\smallskip}

\newcommand{\kfOXItem}{\item}

% v16: O/X 그래픽 박스 키움 (18×14 → 24×20pt) + 영역색 강조
\newcommand{\kfOX}{%
  \hfill\nolinebreak
  \tikz[baseline=-0.9ex]{%
    \node[draw=kfPrimary, fill=kfPrimary!5, line width=0.6pt, rounded corners=3pt,
          minimum width=24pt, minimum height=20pt, inner sep=0pt,
          font=\normalsize\bfseries\color{kfPrimary}] (ob) {O};%
    \node[draw=kfMute, fill=kfMute!5, line width=0.6pt, rounded corners=3pt,
          minimum width=24pt, minimum height=20pt, inner sep=0pt,
          font=\normalsize\bfseries\color{kfMute},
          right=4pt of ob] (xb) {X};%
  }%
}

% ============================================================
% activity-connect.tex
% 좌우 선 긋기 활동 (2단 양식 + 중앙 여백)
% 사용:
%   \begin{kfConnect}
%     \kfPair{사르트르}{실존은 본질에 앞선다}
%     \kfPair{니체}{초인 사상}
%     \kfPair{하이데거}{현존재}
%   \end{kfConnect}
% ============================================================

\newenvironment{kfConnect}%
  {\par\smallskip%
   \renewcommand{\arraystretch}{1.6}%
   \begin{tabular}{@{}>{\raggedleft\arraybackslash}m{4.5cm}
                     @{\hspace{2.5cm}}
                     >{\raggedright\arraybackslash}m{6.5cm}@{}}}%
  {\end{tabular}\par\smallskip}

\newcommand{\kfPair}[2]{%
  \tikz[baseline]{\node[draw=kfRule, fill=kfPrimaryLight!30, rounded corners=2pt,
        inner sep=4pt, font=\small]{#1};} \tikz[baseline]{\node[circle, draw=kfMute, fill=kfPaper, inner sep=1.5pt]{};}
  &
  \tikz[baseline]{\node[circle, draw=kfMute, fill=kfPaper, inner sep=1.5pt]{};} \tikz[baseline]{\node[draw=kfRule, fill=kfNoteBg, rounded corners=2pt,
        inner sep=4pt, font=\small]{#2};} \\
}

% ============================================================
% activity-writing.tex
% 글쓰기 활동 (큰 노트 영역)
% 사용:
%   \kfWriting{독후감}{12}     % 라벨 + 12줄 큰 노트
%   \kfWritingPrompt{프롬프트 텍스트}
% ============================================================

\newcommand{\kfWritingPrompt}[1]{%
  \par\smallskip%
  \begin{tcolorbox}[
    enhanced, breakable=false,
    colback=kfPrimaryLight!30, colframe=kfPrimary,
    boxrule=0.4pt, arc=3pt,
    left=\kfBoxPad, right=\kfBoxPad, top=\kfBoxPad, bottom=\kfBoxPad,
  ]%
  \small\textbf{\color{kfPrimary}글쓰기 안내}\\[2pt]%
  #1
  \end{tcolorbox}\par\smallskip%
}

\newcommand{\kfWriting}[2]{%
  \par\smallskip\needspace{\numexpr#2+2\relax\baselineskip}%
  \begin{tcolorbox}[
    enhanced, breakable=false,
    colback=kfPaper, colframe=kfRule,
    boxrule=0.4pt, arc=2pt,
    left=\kfBoxPad, right=\kfBoxPad, top=\kfBoxPad, bottom=\kfBoxPad,
    title={\footnotesize\bfseries\color{kfMute} #1},
    coltitle=kfMute, colbacktitle=kfPaper,
    attach boxed title to top left={yshift=-1.5mm, xshift=6pt},
    boxed title style={colframe=kfRule, boxrule=0.3pt, arc=1pt}
  ]%
  \foreach \i in {1,...,#2}{%
    \textcolor{kfRule}{\rule{\linewidth}{0.3pt}}\\[7pt]%
  }%
  \end{tcolorbox}\par\smallskip%
}

% ============================================================
% activity-structure.tex
% 구조도 (트리 + 빈칸)
% 사용:
%   \begin{kfStructure}
%     \kfNode{0}{서론}{문제 제기}
%     \kfNode{1}{본론}{(\,\quad\,)}
%     \kfNode{1}{본론}{근거 2}
%     \kfNode{0}{결론}{주장 강화}
%   \end{kfStructure}
% ============================================================

\newenvironment{kfStructure}%
  {\par\smallskip\needspace{6\baselineskip}%
   \begin{tcolorbox}[
     enhanced, breakable=false,
     colback=kfPaper, colframe=kfRule,
     boxrule=0.4pt, arc=2pt,
     left=\kfBoxPad, right=\kfBoxPad, top=\kfBoxPad, bottom=\kfBoxPad,
     title={\footnotesize\bfseries\color{kfMute} 글의 구조},
     coltitle=kfMute, colbacktitle=kfPrimaryLight!40,
     attach boxed title to top left={yshift=-1.5mm, xshift=6pt},
     boxed title style={colframe=kfRule, boxrule=0.3pt, arc=1pt}
   ]%
   \begin{tabbing}
     \hspace{1.4cm}\=\hspace{1.4cm}\=\hspace{8cm}\kill}%
  {\end{tabbing}\end{tcolorbox}\par\smallskip}

% 노드: #1=들여쓰기 단계 (0~3), #2=라벨, #3=내용
\newcommand{\kfNode}[3]{%
  \ifcase#1\relax%
    \>\>\textbf{\color{kfPrimary} ▶ #2}\quad #3\\
  \or
    \>\>\quad\textbf{\color{kfMute} ◦ #2}\quad #3\\
  \or
    \>\>\quad\quad\textbf{\color{kfMute} - #2}\quad #3\\
  \or
    \>\>\quad\quad\quad\textbf{\color{kfMute} \textperiodcentered\ #2}\quad #3\\
  \fi
}

% 빈칸 노드
\newcommand{\kfNodeBlank}[2]{\kfNode{#1}{#2}{(\hspace{2.5cm})}}

% ============================================================
% activity-sentence-reading.tex
% 문장 독해 활동 — 문장 박스 1개 + 그에 묶인 문제 N개를 한 minipage로
% 사용:
%   \begin{kfSentenceUnit}{1}{"바람은 내 귀에 속삭이며..."}
%     \kfSentQ{(1)}{서술어 '속삭이며'의 주어는?}
%     \begin{kfChoices}
%       \kfChoice 바람은
%       ...
%     \end{kfChoices}
%   \end{kfSentenceUnit}
% ============================================================

% kfSentenceBox 매크로는 passage-note.tex에 정의됨

\newenvironment{kfSentenceUnit}[2]%
  {\par\smallskip\needspace{6\baselineskip}%
   \noindent\begin{minipage}{\linewidth}%
   \samepage%
   \kfSentenceBox{문장 #1}{#2}%
   \par\smallskip%
  }%
  {\end{minipage}\par\smallskip}

% 같은 문장에 묶인 소문항 (문장 안내 + 발문)
\newcommand{\kfSentQ}[2]{%
  \par\noindent\textbf{#1}~#2\par\smallskip%
}

% 헤더 없이 문장만 있는 묶음 (sentence_ref가 없는 일반 문항)
\newenvironment{kfSentencelessUnit}%
  {\par\smallskip\noindent\begin{minipage}{\linewidth}\samepage}%
  {\end{minipage}\par\smallskip}

% ============================================================
% chapter-cover.tex (v6)
% 챕터 진입 페이지 (표지) — 하단에 영역 목차 추가
% 사용:
%   \kfChapterCover{1}{근대성과 자아}{Chapter 1}{이번 챕터에서는 ...}
%   \begin{kfChapterAreaList}
%     \kfChapterAreaItem{어휘}{kfAreaVocab}{핵심 어휘 12개}
%     ...
%   \end{kfChapterAreaList}
%
%   영역 목차가 없을 때는 \kfChapterCoverEnd 호출
% ============================================================

\newcommand{\kfChapterCover}[4]{%
  \clearpage%
  \thispagestyle{kfBlank}%
  \kfMarkChapter{Chapter #1. #2}%
  \kfChapterCoverBg%
  % v18.5: 챕터 표지 우상단에 로고
  \begin{tikzpicture}[remember picture, overlay]
    \node[anchor=north east, inner sep=0pt] at ($(current page.north east)+(-18mm,-18mm)$) {\kfLogoMd};
  \end{tikzpicture}%
  \vspace*{20mm}%
  \noindent
  \begin{tikzpicture}[remember picture, overlay]
    % 좌측 액센트 바
    \fill[kfPrimary] (current page.west) ++(12mm,25mm) rectangle ++(5mm, -50mm);
    % 챕터 번호 배지 (이중 원, 약간 작게)
    \node[anchor=north west, inner sep=0pt] at ($(current page.north west)+(24mm,-45mm)$) {%
      \begin{tikzpicture}
        \draw[kfPrimary, line width=1pt] (0,0) circle (10mm);
        \draw[kfPrimary, line width=0.4pt] (0,0) circle (8mm);
        \node[font=\normalsize\bfseries\color{kfPrimary}] at (0,2mm) {CHAPTER};
        \node[font=\LARGE\bfseries\color{kfPrimary}] at (0,-3mm) {#1};
      \end{tikzpicture}%
    };
  \end{tikzpicture}%
  \vspace{42mm}%
  \par\noindent\hspace{32mm}%
  \begin{minipage}{0.65\linewidth}%
    {\footnotesize\color{kfMute} #3}\\[4pt]
    {\LARGE\bfseries\color{kfPrimary} #2}\\[8pt]
    {\color{kfRule}\rule{50mm}{0.5pt}}\\[6pt]
    {\small\color{kfInk} #4}
  \end{minipage}%
  \par\vspace{14mm}%
}

% v6 / v18.5 / v18.6: 챕터 표지의 영역 목차 — 흰색 반투명 박스
% v18.6: 배경 가로선/도형과 안 겹치도록 TikZ overlay로 우하단 절대 위치
\makeatletter
% 영역 항목들을 모아 둘 박스
\newsavebox{\kf@arealistbox}
\newenvironment{kfChapterAreaList}{%
  \begin{lrbox}{\kf@arealistbox}%
  \begin{minipage}{0.62\linewidth}%
    \begin{tcolorbox}[%
      enhanced, colback=white, colframe=white,
      opacityback=0.92, opacityframe=0,
      boxrule=0pt, arc=4pt,
      width=\linewidth,
      top=8pt, bottom=8pt, left=10pt, right=10pt,
    ]%
      {\footnotesize\bfseries\color{kfMute} 이 챕터의 영역}\par\vspace{4pt}%
      {\color{kfRule}\hrule height 0.3pt}\par\vspace{4pt}%
      \begin{itemize}[leftmargin=14pt, itemsep=2pt, topsep=0pt, label={\color{kfPrimary!70}\textbullet}]%
}{%
      \end{itemize}%
    \end{tcolorbox}%
  \end{minipage}%
  \end{lrbox}%
  % v18.7: 배경 상단 가로선 ~ 하단 가로선 사이의 중간 영역에 배치
  % 가로선이 내용을 가리지 않도록 박스 중심이 두 선 사이 중점에 위치
  \begin{tikzpicture}[remember picture, overlay]%
    \node[anchor=center, inner sep=0pt]
      at ($(current page.south)+(0mm,90mm)$)
      {\usebox{\kf@arealistbox}};%
  \end{tikzpicture}%
  \vfill%
  \noindent\hfill\kfSeriesBadge\hspace{12mm}%
  \clearpage%
}
\makeatother

% 영역 항목: \kfChapterAreaItem{영역명}{kfArea색}{부제}
\makeatletter
\newcommand{\kfChapterAreaItem}[3]{%
  \item {\small\bfseries\color{#2} #1}%
  \def\kf@cci@sub{#3}%
  \ifx\kf@cci@sub\@empty\else
    {\small\color{kfMute}\, \textemdash\, #3}%
  \fi
}
\makeatother

% 영역 목록 없이 cover만 (legacy)
\newcommand{\kfChapterCoverEnd}{%
  \vfill%
  \noindent\hfill\kfSeriesBadge\hspace{12mm}%
  \clearpage%
}

% ============================================================
% area-header.tex (v16)
% 영역 시작 표제 — 시리즈별 차별화 라벨 + 큰 영역명 + 부제
% PAGE_BREAK_POLICY.md §1.4
%
% v16 (디자인 고도화 Phase 1):
%   - 시리즈별 라벨 박스 추가:
%     소쉬르 = AREA START (영역색 채움 박스)
%     프레게 = SECTION OPEN (영역색 테두리)
%     러셀   = RUSSELL (영역색 라이트 박스)
%     비트   = WITTGENSTEIN (회색 라이트 박스)
%   - 라벨 위치: 영역명 위 좌측
% ============================================================

\makeatletter

% 시리즈 라벨 텍스트
\newcommand{\kf@serieslabel}{%
  \ifcase\kf@series\relax
    AREA START%
  \or
    AREA START%
  \or
    SECTION OPEN%
  \or
    RUSSELL%
  \or
    WITTGENSTEIN%
  \fi
}

% 시리즈 라벨 박스 스타일 (#1 = 영역색)
\newcommand{\kf@labelbox}[1]{%
  \ifcase\kf@series\relax
    % 0/1 = 소쉬르: 영역색 채움
    \tikz[baseline=-2pt]{\node[fill=#1, text=white,
      rounded corners=3pt, inner xsep=6pt, inner ysep=2pt,
      font=\scriptsize\bfseries]{\kf@serieslabel};}%
  \or
    % 1 = 소쉬르: 영역색 채움
    \tikz[baseline=-2pt]{\node[fill=#1, text=white,
      rounded corners=3pt, inner xsep=6pt, inner ysep=2pt,
      font=\scriptsize\bfseries]{\kf@serieslabel};}%
  \or
    % 2 = 프레게: 영역색 테두리
    \tikz[baseline=-2pt]{\node[draw=#1, line width=0.6pt, text=#1,
      rounded corners=3pt, inner xsep=6pt, inner ysep=2pt,
      font=\scriptsize\bfseries]{\kf@serieslabel};}%
  \or
    % 3 = 러셀: 영역색 라이트 박스
    \tikz[baseline=-2pt]{\node[fill=#1!12, text=#1,
      rounded corners=2pt, inner xsep=5pt, inner ysep=1.5pt,
      font=\scriptsize\bfseries]{\kf@serieslabel};}%
  \or
    % 4 = 비트: 회색 미니 박스
    \tikz[baseline=-2pt]{\node[fill=kfMute!10, text=kfMute,
      rounded corners=2pt, inner xsep=5pt, inner ysep=1.5pt,
      font=\scriptsize\bfseries]{\kf@serieslabel};}%
  \fi
}

% v17: 시리즈+영역별 Canva 배경 자동 매핑
% \kf@hasbg=1로 set되면 \kfAreaStart에서 라벨 박스 출력 생략
\def\kf@areaVocab{어휘}\def\kf@areaGrammar{문법}\def\kf@areaConcept{개념}
\def\kf@areaLit{문학}\def\kf@areaNonlit{비문학}
\newif\ifkf@bgon
% 파일 있으면 배경 적용 + boolean true. 없으면 nothing.
\newcommand{\kfBgSet}[1]{%
  \IfFileExists{#1.pdf}{\kfPageBg{#1}\kf@bgontrue}{}%
}
\newcommand{\kf@trybg}[1]{%
  % v18.4: 영역 배경 PDF 비활성화 — 큰 도형이 본문 영역을 침범해
  % "영역 헤더 후 빈 페이지" 문제를 일으킴. 라벨 박스로만 영역 표시.
  \kf@bgonfalse%
}

% Note: 러셀(rus_*)/비트(wit_*) 배경 PDF 미존재 시 \kfPageBg가 에러 → \IfFileExists로 가드
% (현재는 파일이 모두 존재한다고 가정. 부분 제공 시 cls 수정 필요)

\newcommand{\kfAreaStart}[3]{%
  \clearpage%
  \thispagestyle{fancy}%
  \kfMarkArea{#1}%
  \kf@trybg{#1}%
  \vspace*{1mm}%
  % 배경 없으면 라텍스 라벨 박스 출력 (배경 있으면 일러스트가 라벨 역할)
  \ifkf@bgon
    \vspace*{4mm}%
  \else
    \noindent\kf@labelbox{#2}\par\vspace{4pt}%
  \fi
  % 영역명 + 부제
  \noindent
  \begin{tikzpicture}
    \node[anchor=west, inner sep=0pt] (bar) at (0,0) {\color{#2}\rule{3pt}{22pt}};
    \node[anchor=base west, inner sep=0pt, xshift=8pt, yshift=2pt] (lbl) at (bar.east |- 0,0) {%
      {\LARGE\bfseries\color{#2} #1}%
    };
    \def\kf@areasub{#3}%
    \ifx\kf@areasub\@empty\else
      \node[anchor=base west, inner sep=0pt, xshift=8pt, font=\normalsize\color{kfMute}]
        at (lbl.base east) {\,\textperiodcentered\, #3};%
    \fi
  \end{tikzpicture}%
  \par\vspace{2pt}
  {\color{#2!70}\hrule height 1pt}%
  \par\vspace{1pt}%
  {\color{kfRule}\hrule height 0.3pt}%
  \par\vspace{4pt}%
  \needspace{6\baselineskip}\nopagebreak[4]%
  \global\kf@areaJustOpenedtrue% v18.5: 첫 컨텐츠 needspace 무시
}
\makeatother

% 시리즈/영역 단축 매크로
\newcommand{\kfStartVocab}[1]{\kfAreaStart{어휘}{kfAreaVocab}{#1}}
\newcommand{\kfStartGrammar}[1]{\kfAreaStart{문법}{kfAreaGrammar}{#1}}
\newcommand{\kfStartConcept}[1]{\kfAreaStart{개념}{kfAreaConcept}{#1}}
\newcommand{\kfStartLit}[1]{\kfAreaStart{문학}{kfAreaLiterature}{#1}}
\newcommand{\kfStartNonlit}[1]{\kfAreaStart{비문학}{kfAreaNonlit}{#1}}
\newcommand{\kfStartWeekly}[1]{\kfAreaStart{실력 확인}{kfAreaWeekly}{#1}}

% ============================================================
% section-header.tex
% 섹션 시작 라벨 헤더 — [지문] / [활동] / [문제] 라벨
% 좌측 액센트 바 + 라벨 + 부제 (영역 액센트 색)
%
% 사용:
%   \kfSecHeader{kfAreaLiterature}{지문}{빼앗긴 들에도 봄은 오는가 / 이상화}
%   \kfSecHeader{kfAreaConcept}{활동}{내용 확인}
%   \kfSecHeader{kfAreaNonlit}{문제}{}
%
% 헤더 직후 페이지 분할 금지: \needspace{4\baselineskip} + \nopagebreak
% ============================================================

\providecommand{\kfSecHeader}[3]{%
  \par\needspace{10\baselineskip}\filbreak%
  \vspace{3pt}%
  \noindent
  \begin{tikzpicture}
    \node[anchor=west, inner sep=0pt] (bar) at (0,0) {\color{#1}\rule{3pt}{14pt}};
    \node[anchor=west, inner sep=0pt, xshift=6pt, font=\normalsize\bfseries\color{#1}]
       (lbl) at (bar.east) {[\,#2\,]};
    \node[anchor=west, inner sep=0pt, xshift=4pt, font=\small\color{kfMute}]
       at (lbl.east) {#3};
  \end{tikzpicture}%
  \par\vspace{1pt}%
  {\color{#1!50}\hrule height 0.4pt}%
  \par\vspace{2pt}\nopagebreak%
}

% 영역별 단축 매크로
\newcommand{\kfSecPassageVocab}[1]{\kfSecHeader{kfAreaVocab}{지문}{#1}}
\newcommand{\kfSecPassageGrammar}[1]{\kfSecHeader{kfAreaGrammar}{지문}{#1}}
\newcommand{\kfSecPassageConcept}[1]{\kfSecHeader{kfAreaConcept}{지문}{#1}}
\newcommand{\kfSecPassageLit}[1]{\kfSecHeader{kfAreaLiterature}{지문}{#1}}
\newcommand{\kfSecPassageNonlit}[1]{\kfSecHeader{kfAreaNonlit}{지문}{#1}}
\newcommand{\kfSecPassageWeekly}[1]{\kfSecHeader{kfAreaWeekly}{지문}{#1}}

\newcommand{\kfSecActVocab}[1]{\kfSecHeader{kfAreaVocab}{활동}{#1}}
\newcommand{\kfSecActGrammar}[1]{\kfSecHeader{kfAreaGrammar}{활동}{#1}}
\newcommand{\kfSecActConcept}[1]{\kfSecHeader{kfAreaConcept}{활동}{#1}}
\newcommand{\kfSecActLit}[1]{\kfSecHeader{kfAreaLiterature}{활동}{#1}}
\newcommand{\kfSecActNonlit}[1]{\kfSecHeader{kfAreaNonlit}{활동}{#1}}
\newcommand{\kfSecActWeekly}[1]{\kfSecHeader{kfAreaWeekly}{활동}{#1}}

\newcommand{\kfSecQVocab}[1]{\kfSecHeader{kfAreaVocab}{문제}{#1}}
\newcommand{\kfSecQGrammar}[1]{\kfSecHeader{kfAreaGrammar}{문제}{#1}}
\newcommand{\kfSecQConcept}[1]{\kfSecHeader{kfAreaConcept}{문제}{#1}}
\newcommand{\kfSecQLit}[1]{\kfSecHeader{kfAreaLiterature}{문제}{#1}}
\newcommand{\kfSecQNonlit}[1]{\kfSecHeader{kfAreaNonlit}{문제}{#1}}
\newcommand{\kfSecQWeekly}[1]{\kfSecHeader{kfAreaWeekly}{문제}{#1}}


\begin{document}

\thispagestyle{kfBlank}
\kfBookCoverBg
\vspace*{\kfBookCoverVspace}
\begin{center}
{\kfLogoLg}\\[14pt]
{\fontsize{72}{84}\selectfont\bfseries\kfBookCoverColor 프레게}\\[18pt]
{\fontsize{28}{32}\selectfont\kfBookCoverSubColor \textsc{Level} \textbf{1}}\\[22pt]
{\large\kfBookCoverSubColor 제 1 권 \quad\textbar\quad 챕터 1\textendash{} 4}
\end{center}
\vfill
\hfill\kfSeriesBadge\hspace{15mm}
\clearpage
\kfChapterCover{1}{프레게1 (챕터1)}{Chapter 1}{이번 챕터의 학습 내용을 시작합니다.}
\begin{kfChapterAreaList}
\kfChapterAreaItem{어휘}{kfAreaVocab}{}
\kfChapterAreaItem{문법}{kfAreaGrammar}{중심 생각과 뒷받침 문장}
\kfChapterAreaItem{개념}{kfAreaConcept}{생각을 정리하는 것이 학습에 미치는 영향}
\kfChapterAreaItem{문학}{kfAreaLiterature}{「지우는 요즘 고민이 많았다. 다음 주에 …」}
\kfChapterAreaItem{비문학}{kfAreaNonlit}{생각을 분류하고 조직하는 법}
\kfChapterAreaItem{실력확인}{kfAreaWeekly}{}
\end{kfChapterAreaList}

\kfStartVocab{어휘 영역 학습}


\kfActivityBg \par\kfMaybeNeedspace{24\baselineskip}\noindent{\kfTipMark\,\,}{\small\color{kfInk} [1-4] 다음 빈칸에 알맞은 어휘를 넣으세요.}\par\nopagebreak[4]\smallskip\nopagebreak[4]
\begin{kfBlankList}
\kfBlankItem 도서관에서는 책을 주제별로 [　　　　　]하여 찾기 쉽게 만든다. (종류에 따라 나누어 구분하는 어휘를 넣으세요)
\kfBlankItem 생각이 [　　　　　]에 빠지면 좋은 아이디어가 떠오르지 않는다. (뒤죽박죽 섞여 어지러운 상태를 나타내는 어휘를 넣으세요)
\kfBlankItem 공부를 [　　　　　]으로 하면 더 효과적으로 학습할 수 있다. (일정한 틀에 따라 짜임새가 있다는 의미의 어휘를 넣으세요)
\kfBlankItem 같은 [　　　　　]에 속하는 동물들은 비슷한 특징을 가지고 있다. (같은 성질을 가진 것들의 묶음을 나타내는 어휘를 넣으세요)
\end{kfBlankList}

\kfSecHeader{kfAreaVocab}{문제}{}

\begin{kfQBlock}{01}{객관식}
\kfQStem{'정리'의 반의어로 가장 적절한 것은?}
\begin{kfChoices}
\kfChoice 분류
\kfChoice 혼란
\kfChoice 배열
\kfChoice 조직
\end{kfChoices}
\end{kfQBlock}
\begin{kfQBlock}{02}{객관식}
\kfQStem{다음 중 '체계'와 가장 가까운 의미를 가진 것은?}
\begin{kfChoices}
\kfChoice 질서
\kfChoice 혼란
\kfChoice 방해
\kfChoice 실수
\end{kfChoices}
\end{kfQBlock}
\begin{kfQBlock}{03}{객관식}
\kfQStem{다음 문장에서 밑줄 친 부분과 바꿔 쓸 수 있는 말은?}
\begin{kfEvidence}선생님께서는 복잡한 내용을 \uline{체계적으로} 설명해 주셨다.\end{kfEvidence}
\begin{kfChoices}
\kfChoice 어지럽게
\kfChoice 짜임새 있게
\kfChoice 대충대충
\kfChoice 느리게
\end{kfChoices}
\end{kfQBlock}
\begin{kfQBlock}{04}{객관식}
\kfQStem{'구분'과 '분류'의 차이를 설명한 것으로 가장 적절한 것은?}
\begin{kfChoices}
\kfChoice 구분은 차이점을 중심으로, 분류는 공통점을 중심으로 나눈다.
\kfChoice 구분과 분류는 완전히 같은 의미이다.
\kfChoice 구분은 큰 것을 작게, 분류는 작은 것을 크게 만든다.
\kfChoice 구분은 시간순으로, 분류는 공간순으로 나눈다.
\end{kfChoices}
\end{kfQBlock}
\begin{kfQBlock}{05}{서술형}
\kfQStem{'배열'을 사용한 문장}
\kfAnswerNote{답안}{4}
\end{kfQBlock}
\begin{kfQBlock}{06}{서술형}
\kfQStem{'질서'를 사용한 문장}
\kfAnswerNote{답안}{4}
\end{kfQBlock}


\kfStartGrammar{중심 생각과 뒷받침 문장}


\kfPassageNoteNT{%
중심 생각(주제문): 문단에서 가장 중요한 내용을 담고 있는 문장.

뒷받침 문장(보조문): 중심 생각을 설명하거나 예시를 들어 구체화하는 문장들.
}

\kfActivityBg \kfSecHeader{kfAreaGrammar}{활동}{지시어 연결}
\par\kfMaybeNeedspace{18\baselineskip}\noindent{\kfTipMark\,\,}{\small\color{kfInk} 지시어가 가리키는 대상을 찾아 답해 보세요.}\par\nopagebreak[4]\smallskip\nopagebreak[4]
\begin{enumerate}[leftmargin=22pt, itemsep=6pt, topsep=4pt, label=\textbf{\arabic*.}]
\item 다음 문장에서 '이것'과 '이러한 능력'이 가리키는 것을 찾아 쓰세요.
\begin{kfEvidence}"생각을 정리하는 것은 중요합니다. \uline{이것}은 우리의 사고력을 향상시킵니다. 또한 \uline{이러한 능력}은 학습 효율을 높이는 데도 도움이 됩니다."\end{kfEvidence}
\end{enumerate}

\kfActivityBg \kfSecHeader{kfAreaGrammar}{활동}{어법 훈련}
\par\kfMaybeNeedspace{24\baselineskip}\noindent{\kfTipMark\,\,}{\small\color{kfInk} [1-3] 문장의 뜻에 알맞은 낱말을 골라 O표 하세요.}\par\nopagebreak[4]\smallskip\nopagebreak[4]
\begin{enumerate}[leftmargin=22pt, itemsep=6pt, topsep=4pt, label=\textbf{\arabic*.}]
\item 나와 친구는 서로 생김새가 ( 달라요 / 틀려요 ).
\item 선생님께서 모르는 문제의 풀이 방법을 친절하게 ( 가르쳐 / 가리켜 ) 주셨다.
\item 숙제 가져오는 것을 깜빡 ( 잃어버렸다 / 잊어버렸다 ).
\end{enumerate}

\kfSecHeader{kfAreaGrammar}{문제}{}

\begin{kfQBlock}{01}{객관식}
\kfQStem{<보기>에서 중심 생각을 담은 문장은?}
\begin{kfEvidence}㉠ 규칙적인 운동은 건강을 유지하는 가장 좋은 방법이다. \\\relax
㉡ 운동을 하면 심장이 튼튼해진다. \\\relax
㉢ 또한 근육이 강해지고 뼈가 단단해진다. \\\relax
㉣ 무엇보다 스트레스가 줄어들어 정신 건강에도 도움이 된다.\end{kfEvidence}
\begin{kfChoices}
\kfChoice ㉠ 문장
\kfChoice ㉡ 문장
\kfChoice ㉢ 문장
\kfChoice ㉣ 문장
\end{kfChoices}
\end{kfQBlock}
\begin{kfQBlock}{02}{객관식}
\kfQStem{다음 뒷받침 문장들을 모두 아우르는 '중심 문장'으로 가장 알맞은 것은 무엇인가요?}
\begin{kfEvidence}아침에는 시원한 물 한 잔으로 시작한다. \\\relax
식사할 때마다 물을 마신다. \\\relax
운동 후에는 충분한 수분을 보충한다.\end{kfEvidence}
\begin{kfChoices}
\kfChoice 아침에 일찍 일어나는 습관을 갖자.
\kfChoice 운동을 규칙적으로 해야 튼튼해진다.
\kfChoice 식사를 거르지 말고 규칙적으로 하자.
\kfChoice 물을 자주 마시는 습관을 기르자.
\end{kfChoices}
\end{kfQBlock}


\kfStartConcept{생각을 정리하는 것이 학습에 미치는 영향}


\kfPassageNoteNT{%
　우리가 배우는 모든 정보는 마치 서랍 안에 넣는 물건과 같습니다. 서랍을 체계적으로 정리하면 필요한 물건을 쉽게 찾을 수 있듯이, 생각을 분류하고 조직하면 지식을 효율적으로 머릿속에 담을 수 있습니다. 

　생각을 정리하는 능력이 뛰어나면 단순히 숙제를 잘하는 것을 넘어, 새로운 지식을 더 빨리 이해하고, 복잡한 문제의 핵심을 빠르게 파악할 수 있습니다. 예를 들어, 수학 문제를 풀 때 어느 범주에 속하는지, 어떤 단계부터 풀어야 할지 우선순위를 정하는 과정 모두 생각 정리 능력에 포함됩니다. 정리된 생각이 바로 효율적인 학습의 시작점입니다.
\par\medskip\begin{itemize}[leftmargin=18pt, itemsep=2pt, topsep=2pt, label=\textbullet]
\item 분류 — 과목별·주제별 노트 나누기
\item 조직 — 개요·마인드맵으로 구조화
\item 우선순위 — 수학 문제를 단계별 순서로 풀기
\end{itemize}
}

\kfActivityBg \par\kfMaybeNeedspace{24\baselineskip}\noindent{\kfTipMark\,\,}{\small\color{kfInk} [1-3] 위 글의 내용과 일치하면 O, 일치하지 않으면 X 표시를 하세요.}\par\nopagebreak[4]\smallskip\nopagebreak[4]
\begin{kfOXList}
\kfOXItem 생각을 체계적으로 정리하면 지식을 머릿속에 효율적으로 담을 수 있다. ( ○ /  ) \kfOX
\kfOXItem 생각 정리 능력은 오직 수학 문제를 풀 때만 도움이 된다. ( ○ /  ) \kfOX
\kfOXItem 정리된 생각은 복잡한 문제의 핵심을 빠르게 파악하도록 도와준다. ( ○ /  ) \kfOX
\end{kfOXList}


\kfStartLit{「지우는 요즘 고민이 많았다. 다음 주에 …」}


\kfPassageNote{}{%
　지우는 요즘 고민이 많았다. 다음 주에 있을 발표 준비도 해야 하고, 엄마가 부탁한 심부름도 까먹으면 안 되고, 친구와의 약속도 지켜야 했다. 게다가 어제 본 재미있는 영화 이야기가 자꾸 떠올라서 집중이 되지 않았다.

"아, 머리가 복잡해!" 지우는 책상에 엎드려 한숨을 쉬었다. 그때 옆자리에 앉은 수민이가 말했다. "지우야, 네 책상 좀 봐. 책상이 그렇게 어지러우니까 머릿속도 복잡한 거 아닐까?"

　지우는 자신의 책상을 둘러보았다. 교과서, 노트, 필통, 간식 봉지까지 온갖 물건들이 뒤섞여 있었다. "그런데 이게 내 머리랑 무슨 상관이야?" 지우가 물었다.

　수민이는 자신의 깔끔한 책상을 가리키며 설명했다. "나는 매일 아침 책상을 정리하면서 오늘 할 일도 함께 정리해. 교과서는 교과서끼리, 노트는 노트끼리 모아두고, 오늘 할 일은 메모지에 적어서 잘 보이는 곳에 붙여둬. 그러면 머릿속도 깔끔해지는 기분이야."

　지우는 반신반의하면서도 수민이의 방법을 따라해 보기로 했다. 먼저 책상 위의 물건들을 종류별로 나누었다. 그다음 오늘 할 일을 메모지에 하나씩 적었다. 

　'1. 발표 자료 찾기, 2. 우유 사오기, 3. 민호와 4시에 도서관에서 만나기.'

　신기하게도 책상이 정리되자 머릿속도 한결 가벼워지는 느낌이었다. 지우는 이제 무엇부터 해야 할지 명확히 알 수 있었다. "와, 정말 신기하다! 생각을 정리하는 것도 연습이 필요한 거구나." 지우는 앞으로 매일 아침 책상 정리와 함께 하루를 시작하기로 마음먹었다.
}


\kfActivityBg \par\kfMaybeNeedspace{18\baselineskip}\noindent{\kfTipMark\,\,}{\small\color{kfInk} [1-1] 창의 활동: 나만의 정리 방법 만들기}\par\nopagebreak[4]\smallskip\nopagebreak[4]
\kfWritingPrompt{지우처럼 여러분도 머릿속이 복잡할 때가 있을 거예요. 자신만의 생각 정리 방법을 만들어 보세요.

예시: "나는 '신호등 정리법'을 사용해요. 빨간색은 오늘 꼭 해야 할 일, 노란색은 이번 주 안에 해야 할 일, 초록색은 시간이 있을 때 하면 되는 일로 구분해요."

나의 생각 정리 방법:}
\kfWriting{답안 작성}{8}

\kfActivityBg \par\kfMaybeNeedspace{18\baselineskip}\noindent{\kfTipMark\,\,}{\small\color{kfInk} 다음 조건에 맞게 글을 써 보세요.}\par\nopagebreak[4]\smallskip\nopagebreak[4]
\kfWritingPrompt{생각을 정리하면 어떤 점이 좋을까요? 지우의 경험을 참고하여, 정리가 주는 긍정적인 효과를 두 가지 이상 써 보세요.}
\kfWriting{답안 작성}{8}

\kfSecHeader{kfAreaLiterature}{문제}{}

\begin{kfQBlock}{01}{객관식}
\kfQStem{지우가 고민하는 일로 적절하지 \uline{\underline{않은}} 것은?}
\begin{kfChoices}
\kfChoice 다음 주 발표 준비
\kfChoice 엄마의 심부름
\kfChoice 친구와의 약속
\kfChoice 수학 시험 공부
\end{kfChoices}
\end{kfQBlock}
\begin{kfQBlock}{02}{객관식}
\kfQStem{지우가 메모지에 적은 할 일의 순서로 올바른 것은?}
\begin{kfChoices}
\kfChoice 우유 사오기 → 발표 자료 찾기 → 민호 만나기
\kfChoice 발표 자료 찾기 → 우유 사오기 → 민호 만나기
\kfChoice 민호 만나기 → 우유 사오기 → 발표 자료 찾기
\kfChoice 발표 자료 찾기 → 민호 만나기 → 우유 사오기
\end{kfChoices}
\end{kfQBlock}
\begin{kfQBlock}{03}{객관식}
\kfQStem{이야기에서 '책상 정리'와 '생각 정리'의 관계를 가장 잘 설명한 것은?}
\begin{kfChoices}
\kfChoice 책상 정리는 머릿속 생각 정리와 관련이 거의 없다.
\kfChoice 책상을 정리하면 생각도 함께 정리되는 효과가 있다.
\kfChoice 생각을 먼저 정리한 후에만 책상을 정리할 수 있다.
\kfChoice 책상이 오히려 어지러울수록 창의적인 생각이 떠오른다.
\end{kfChoices}
\end{kfQBlock}
\begin{kfQBlock}{04}{객관식}
\kfQStem{이야기에서 지우의 심리 변화 과정을 올바르게 나타낸 것은?}
\begin{kfChoices}
\kfChoice 자신만만함 → 혼란스러움 → 후회
\kfChoice 혼란스러움 → 의구심 → 깨달음
\kfChoice 즐거움 → 슬픔 → 분노
\kfChoice 무관심 → 호기심 → 실망
\end{kfChoices}
\end{kfQBlock}
\begin{kfQBlock}{05}{객관식}
\kfQStem{수민이의 성격으로 가장 적절한 것은?}
\begin{kfChoices}
\kfChoice 계획적이고 체계적이다
\kfChoice 충동적이고 즉흥적이다
\kfChoice 소극적이고 내성적이다
\kfChoice 비판적이고 냉소적이다
\end{kfChoices}
\end{kfQBlock}
\begin{kfQBlock}{06}{객관식}
\kfQStem{이 이야기의 주제로 가장 적절한 것은?}
\begin{kfChoices}
\kfChoice 친구의 소중함
\kfChoice 책상 청소를 해야 하는 이유
\kfChoice 생각을 체계적으로 정리하는 것의 필요성
\kfChoice 학업 성취의 비결
\end{kfChoices}
\end{kfQBlock}
\begin{kfQBlock}{07}{단답형}
\kfQStem{수민이가 지우에게 제안한 방법은 무엇인가요?}
\kfAnswerLines{1}
\end{kfQBlock}
\begin{kfQBlock}{08}{서술형}
\kfQStem{지우가 이야기 끝에서 깨달은 교훈을 서술하세요.}
\kfAnswerNote{답안}{4}
\end{kfQBlock}


\kfStartNonlit{생각을 분류하고 조직하는 법}


\kfPassageNote{생각을 분류하고 조직하는 법}{%
　우리의 뇌는 하루에도 수많은 생각을 처리한다고 합니다. 이 많은 생각들이 무질서하게 떠다니면 우리는 혼란을 느끼고 집중력이 떨어집니다. 그래서 생각을 체계적으로 분류하고 조직하는 능력이 중요합니다. 생각을 정리하는 것은 단순히 머릿속을 깔끔하게 만드는 것이 아니라, 우리의 사고력과 창의성을 높이는 핵심적인 방법입니다.

　생각을 분류하는 첫 번째 방법은 '범주화'입니다. 범주화란 비슷한 특성을 가진 생각들을 하나의 그룹으로 묶는 것을 말합니다. 예를 들어, 학교 관련 생각, 가족 관련 생각, 취미 관련 생각처럼 주제별로 나눌 수 있습니다. 또한 시간을 기준으로 '오늘 할 일', '이번 주 할 일', '이번 달 할 일'로 분류할 수도 있습니다. 이렇게 범주화를 하면 복잡한 생각들이 단순하고 명확해집니다.

　두 번째 방법은 '우선순위 정하기'입니다. 모든 생각이 똑같이 중요한 것은 아닙니다. 긴급하고 중요한 것, 중요하지만 급하지 않은 것, 급하지만 중요하지 않은 것, 급하지도 중요하지도 않은 것으로 나누어 볼 수 있습니다. 이렇게 우선순위를 정하면 어떤 일부터 처리해야 할지 쉽게 결정할 수 있습니다.

　세 번째 방법은 '시각화'입니다. 생각을 눈에 보이는 형태로 표현하면 더 쉽게 정리할 수 있습니다. 마인드맵, 표, 그래프 등을 활용하여 생각을 시각적으로 나타내면 전체적인 구조를 한눈에 파악할 수 있습니다. 특히 마인드맵은 중심 주제에서 가지를 뻗어 나가며 관련된 생각들을 연결하는 방식으로, 창의적인 사고를 돕는 효과적인 도구입니다.

　이러한 방법들을 꾸준히 연습하면 누구나 생각을 체계적으로 정리하는 능력을 기를 수 있습니다. 정리된 생각은 명확한 의사결정을 가능하게 하고, 효율적인 학습과 업무 수행의 기초가 됩니다. 또한 스트레스를 줄이고 창의적인 아이디어를 발견하는 데도 큰 도움이 됩니다. 오늘부터 하루 10분씩이라도 자신의 생각을 정리하는 시간을 가져보는 것은 어떨까요?
}


\kfActivityBg \par\kfMaybeNeedspace{32\baselineskip}\noindent{\kfTipMark\,\,}{\small\color{kfInk} 다음은 각 문단의 핵심 내용을 정리한 것입니다. 빈칸에 알맞은 말을 채워 넣으세요.}\par\nopagebreak[4]\smallskip\nopagebreak[4]
\begin{kfStructure}
\kfNode{0}{}{}
\kfNode{0}{}{}
\kfNode{0}{}{}
\kfNode{0}{}{}
\kfNode{0}{}{}
\end{kfStructure}

\kfActivityBg \par\kfMaybeNeedspace{18\baselineskip}\noindent{\kfTipMark\,\,}{\small\color{kfInk} [1-2] 오늘 하루 동안 떠오른 생각들을 분류해 보세요.}\par\nopagebreak[4]\smallskip\nopagebreak[4]
\kfWritingPrompt{긴급하고 중요한 일 :}
\kfWriting{답안 작성}{8}
\kfWritingPrompt{중요하지만 급하지 않은 일 :}
\kfWriting{답안 작성}{8}

\kfSecHeader{kfAreaNonlit}{문제}{}

\begin{kfQBlock}{01}{객관식}
\kfQStem{이 글의 중심 내용으로 가장 적절한 것은?}
\begin{kfChoices}
\kfChoice 우리 뇌가 처리하는 생각의 개수
\kfChoice 생각을 체계적으로 정리하는 방법
\kfChoice 마인드맵의 장점과 활용법
\kfChoice 스트레스를 줄이는 방법
\end{kfChoices}
\end{kfQBlock}
\begin{kfQBlock}{02}{객관식}
\kfQStem{글쓴이가 제시한 생각 정리 방법이 \uline{\underline{아닌}} 것은?}
\begin{kfChoices}
\kfChoice 범주화
\kfChoice 우선순위 정하기
\kfChoice 시각화
\kfChoice 암기하기
\end{kfChoices}
\end{kfQBlock}
\begin{kfQBlock}{03}{객관식}
\kfQStem{'범주화'의 예시로 적절하지 \uline{\underline{않은}} 것은?}
\begin{kfChoices}
\kfChoice 학교 숙제를 과목별로 나누기
\kfChoice 일정을 오전/오후로 구분하기
\kfChoice 중요한 일을 빨간색으로 표시하기
\kfChoice 책을 장르별로 분류하기
\end{kfChoices}
\end{kfQBlock}
\begin{kfQBlock}{04}{객관식}
\kfQStem{다음 중 '긴급하고 중요한 일'에 해당하는 것은?}
\begin{kfChoices}
\kfChoice 내일까지 제출해야 하는 과제
\kfChoice 다음 달에 있을 가족 여행 계획
\kfChoice 친구의 생일 선물 고르기
\kfChoice 방학 때 읽을 책 목록 만들기
\end{kfChoices}
\end{kfQBlock}
\begin{kfQBlock}{05}{서술형}
\kfQStem{글쓴이가 생각 정리를 권하는 이유를 두 가지 이상 서술하세요.}
\kfAnswerNote{답안}{4}
\end{kfQBlock}


\kfStartWeekly{실력확인 영역 학습}


\noindent\textit{\small (제한시간: 30분 / 총 15문항)}\par\medskip
\par\noindent\textbf{[4-6] 다음 글을 읽고 물음에 답하세요.}\par\medskip
\kfPassageNote{지문 4}{%
"아, 머리가 복잡해!" 지우는 책상에 엎드려 한숨을 쉬었다. 그때 옆자리에 앉은 수민이가 말했다. "지우야, 네 책상 좀 봐. 책상이 그렇게 어지러우니까 머릿속도 복잡한 거 아닐까?" 지우는 자신의 책상을 둘러보았다. 교과서, 노트, 필통, 간식 봉지까지 온갖 물건들이 뒤섞여 있었다. "그런데 이게 내 머리랑 무슨 상관이야?" 지우가 물었다. 수민이는 자신의 깔끔한 책상을 가리키며 설명했다. "나는 매일 아침 책상을 정리하면서 오늘 할 일도 함께 정리해. 교과서는 교과서끼리, 노트는 노트끼리 모아두고, 오늘 할 일은 메모지에 적어서 잘 보이는 곳에 붙여둬. 그러면 머릿속도 깔끔해지는 기분이야." 지우는 반신반의하면서도 수민이의 방법을 따라해 보기로 했다. 먼저 책상 위의 물건들을 종류별로 나누었다. 그다음 오늘 할 일을 메모지에 하나씩 적었다. 신기하게도 책상이 정리되자 머릿속도 한결 가벼워지는 느낌이었다. 지우는 이제 무엇부터 해야 할지 명확히 알 수 있었다.
}
\par\noindent\textbf{[7-9] 다음 글을 읽고 물음에 답하세요.}\par\medskip
\kfPassageNote{지문 7}{%
우리의 뇌는 하루에도 수많은 생각을 처리한다고 합니다. 이 많은 생각들이 무질서하게 떠다니면 우리는 혼란을 느끼고 집중력이 떨어집니다. 그래서 생각을 체계적으로 분류하고 조직하는 능력이 중요합니다. 생각을 정리하는 첫 번째 방법은 '범주화'입니다. 범주화란 비슷한 특성을 가진 생각들을 하나의 그룹으로 묶는 것을 말합니다. 예를 들어, 학교 관련 생각, 가족 관련 생각처럼 주제별로 나눌 수 있습니다. 두 번째 방법은 '우선순위 정하기'입니다. 모든 생각이 똑같이 중요한 것은 아닙니다. 긴급하고 중요한 것부터 순서를 정하면 어떤 일부터 처리해야 할지 쉽게 결정할 수 있습니다. 세 번째 방법은 '\uline{( ㉠ )}'입니다. 마인드맵이나 표 등을 활용하여 생각을 눈에 보이는 형태로 표현하면 전체적인 구조를 한눈에 파악할 수 있습니다.
}
\begin{kfQBlock}{01}{객관식}
\kfQStem{다음 중 '체계적'의 뜻으로 가장 알맞은 것은?}
\begin{kfChoices}
\kfChoice 복잡하고 이해하기 어려운 상태
\kfChoice 일정한 틀이나 규칙에 따라 짜임새가 있는 상태
\kfChoice 뒤죽박죽 섞여서 정신이 없는 상태
\kfChoice 아무런 계획 없이 마음대로 하는 상태
\end{kfChoices}
\end{kfQBlock}

\begin{kfQBlock}{02}{객관식}
\kfQStem{다음 빈칸에 들어갈 알맞은 말은?}
\begin{kfEvidence}"도서관에서는 책을 주제별로 [　　　　　]하여 찾기 쉽게 만든다."\end{kfEvidence}
\begin{kfChoices}
\kfChoice 혼란
\kfChoice 분류
\kfChoice 방해
\kfChoice 무시
\end{kfChoices}
\end{kfQBlock}

\begin{kfQBlock}{03}{객관식}
\kfQStem{위 글에서 수민이가 지적한, 지우의 머릿속이 복잡한 이유는 무엇인가요?}
\begin{kfChoices}
\kfChoice 친구와 심하게 다투어서
\kfChoice 책상이 어지럽게 뒤섞여 있어서
\kfChoice 숙제가 너무 어려워서
\kfChoice 배가 너무 고파서
\end{kfChoices}
\end{kfQBlock}

\begin{kfQBlock}{04}{객관식}
\kfQStem{수민이가 지우에게 제안한 방법으로 알맞지 \uline{\underline{않은}} 것은?}
\begin{kfChoices}
\kfChoice 책상 위의 물건을 종류별로 모아두기
\kfChoice 오늘 할 일을 메모지에 적기
\kfChoice 할 일을 잘 보이는 곳에 붙여두기
\kfChoice 책상 위의 모든 물건을 무조건 한꺼번에 버려서 비우기
\end{kfChoices}
\end{kfQBlock}

\begin{kfQBlock}{05}{객관식}
\kfQStem{이 글의 중심 내용으로 가장 알맞은 것은?}
\begin{kfChoices}
\kfChoice 어지러운 책상이 학습에 미치는 부정적 영향과 사례
\kfChoice 생각을 체계적으로 정리하는 세 가지 방법
\kfChoice 시각화를 활용해 미술 작품을 잘 그리는 비결
\kfChoice 친구의 잘못된 학습 습관을 바로잡아 주는 대화 방법
\end{kfChoices}
\end{kfQBlock}

\begin{kfQBlock}{06}{객관식}
\kfQStem{위 글에서 설명한 '범주화'의 예시로 적절한 것은?}
\begin{kfChoices}
\kfChoice 할 일을 중요한 순서대로 번호를 매긴다.
\kfChoice 생각을 그림으로 그려서 벽에 붙인다.
\kfChoice '과일'이라는 묶음 안에 사과, 배, 포도를 넣는다.
\kfChoice 급한 일부터 먼저 처리한다.
\end{kfChoices}
\end{kfQBlock}

\begin{kfQBlock}{07}{객관식}
\kfQStem{위 글의 ㉠에 들어갈 말로, 생각을 눈에 보이게 표현하는 방법을 무엇이라고 하나요?}
\begin{kfChoices}
\kfChoice 암기
\kfChoice 상상
\kfChoice 시각화
\kfChoice 단순화
\end{kfChoices}
\end{kfQBlock}

\begin{kfQBlock}{08}{객관식}
\kfQStem{다음 문장에서 '누가(주어)'와 '무엇을 했다(서술어)'를 바르게 찾은 것은?}
\begin{kfEvidence}"지우는 수민이의 조언을 듣고 책상을 깨끗하게 정리했습니다."\end{kfEvidence}
\begin{kfChoices}
\kfChoice 지우는 - 조언을
\kfChoice 수민이의 - 듣고
\kfChoice 책상을 - 정리했습니다
\kfChoice 지우는 - 정리했습니다
\end{kfChoices}
\end{kfQBlock}

\begin{kfQBlock}{09}{객관식}
\kfQStem{중심 생각과 뒷받침 문장에 대한 설명으로 바르지 \uline{\underline{않은}} 것은?}
\begin{kfChoices}
\kfChoice 중심 생각은 문단에서 가장 중요한 내용을 담고 있다.
\kfChoice 뒷받침 문장은 중심 생각을 보충하고 설명하는 역할을 한다.
\kfChoice 하나의 문단에는 중심 생각이 여러 개 들어갈 수 있다.
\kfChoice 중심 문장은 문단의 앞이나 뒤에 주로 위치한다.
\end{kfChoices}
\end{kfQBlock}

\begin{kfQBlock}{10}{객관식}
\kfQStem{다음 문단에서 중심 문장은 무엇인가요?}
\begin{kfEvidence}① 규칙적인 운동은 건강을 지키는 좋은 습관입니다. \\\relax
　② 운동을 하면 근육이 튼튼해집니다. \\\relax
　③ 또한 운동은 스트레스를 줄여주는 효과도 있습니다. \\\relax
　④ 밤에 잠을 잘 자게 도와주기도 합니다.\end{kfEvidence}
\begin{kfChoices}
\kfChoice ①번 문장
\kfChoice ②번 문장
\kfChoice ③번 문장
\kfChoice ④번 문장
\end{kfChoices}
\end{kfQBlock}

\begin{kfQBlock}{11}{객관식}
\kfQStem{다음 중 낱말의 쓰임이 바른 문장은?}
\begin{kfChoices}
\kfChoice 표지판이 길을 가르쳤다.
\kfChoice 너와 나는 성격이 서로 틀려.
\kfChoice 선생님이 칠판을 손으로 가르치셨다.
\kfChoice 나는 학교 가는 길을 잃어버렸다.
\end{kfChoices}
\end{kfQBlock}

\begin{kfQBlock}{12}{서술형}
\kfQStem{다음 설명에 해당하는 낱말을 쓰세요.}
\begin{kfEvidence}"같은 성질을 가진 것들의 묶음\end{kfEvidence}
\kfAnswerNote{답안}{4}
\end{kfQBlock}

\begin{kfQBlock}{13}{서술형}
\kfQStem{지우가 책상을 정리한 후에 느낀 점은 무엇인지 한 문장으로 쓰세요.}
\kfAnswerNote{답안}{4}
\end{kfQBlock}

\begin{kfQBlock}{14}{서술형}
\kfQStem{중심 문장을 도와 내용을 더 자세하게 설명하거나 예를 드는 문장을 무엇이라고 하나요?}
\kfAnswerNote{답안}{4}
\end{kfQBlock}

\begin{kfQBlock}{15}{서술형}
\kfQStem{다음 문장의 괄호 안에 알맞은 말을 쓰세요.}
\begin{kfEvidence}"수학 계산을 잘못해서 답이 [　　　　　]." (셈이나 사실이 맞지 않다는 뜻)\end{kfEvidence}
\kfAnswerNote{답안}{4}
\end{kfQBlock}




\clearpage

\kfChapterCover{2}{프레게1 (챕터2)}{Chapter 2}{이번 챕터의 학습 내용을 시작합니다.}
\begin{kfChapterAreaList}
\kfChapterAreaItem{어휘}{kfAreaVocab}{}
\kfChapterAreaItem{문법}{kfAreaGrammar}{글의 짜임 (처음-가운데-끝)}
\kfChapterAreaItem{개념}{kfAreaConcept}{글쓰기에도 지도가 필요해!}
\kfChapterAreaItem{문학}{kfAreaLiterature}{「민준이는 학교에서 돌아오자마자 블록 상자…」}
\kfChapterAreaItem{비문학}{kfAreaNonlit}{개요를 작성하는 방법과 중요성}
\kfChapterAreaItem{실력확인}{kfAreaWeekly}{}
\end{kfChapterAreaList}

\kfStartVocab{어휘 영역 학습}


\kfActivityBg \par\kfMaybeNeedspace{24\baselineskip}\noindent{\kfTipMark\,\,}{\small\color{kfInk} [1-4] 다음 빈칸에 알맞은 어휘를 넣으세요.}\par\nopagebreak[4]\smallskip\nopagebreak[4]
\begin{kfBlankList}
\kfBlankItem 글을 쓰기 전에 [　　　　　] 를 작성하면 내용이 엉뚱한 곳으로 빠지지 않는다. (글을 쓰기 전에 주요 내용을 간추려 짠 틀을 나타내는 어휘를 넣으세요)
\kfBlankItem 건물을 지으려면 튼튼한 [　　　　　] 를 세우는 것이 중요하다. (건물 따위의 뼈대, 건물 뼈대의 짜임새를 나타내는 어휘를 넣으세요)
\kfBlankItem 도서관에 책들이 보기 좋게 [　　　　　] 되어 있다. (알맞은 자리에 나누어 놓음을 나타내는 어휘를 넣으세요)
\kfBlankItem 이 조립 장난감은 설명서의 [　　　　　] 에 따라 만들어야 한다. (순서나 순번을 나타내는 어휘를 넣으세요)
\end{kfBlankList}

\kfSecHeader{kfAreaVocab}{문제}{}

\begin{kfQBlock}{01}{객관식}
\kfQStem{'짜임'과 가장 비슷한 의미를 가진 단어는?}
\begin{kfChoices}
\kfChoice 혼란
\kfChoice 구조
\kfChoice 분해
\kfChoice 파괴
\end{kfChoices}
\end{kfQBlock}
\begin{kfQBlock}{02}{객관식}
\kfQStem{'완성'의 반의어로 가장 적절한 것은?}
\begin{kfChoices}
\kfChoice 성공
\kfChoice 시작
\kfChoice 미완성
\kfChoice 종료
\end{kfChoices}
\end{kfQBlock}
\begin{kfQBlock}{03}{객관식}
\kfQStem{다음 문장에서 밑줄 친 부분과 바꿔 쓸 수 있는 말은?}
\begin{kfEvidence}"건축가가 집의 \uline{설계도}를 그렸다."\end{kfEvidence}
\begin{kfChoices}
\kfChoice 그림
\kfChoice 도면
\kfChoice 낙서
\kfChoice 지도
\end{kfChoices}
\end{kfQBlock}
\begin{kfQBlock}{04}{객관식}
\kfQStem{다음 문장에서 밑줄 친 부분과 바꿔 쓸 수 있는 말은?}
\begin{kfEvidence}"우리는 글의 전체적인 \uline{짜임}을 바꾸기로 했다."\end{kfEvidence}
\begin{kfChoices}
\kfChoice 분량
\kfChoice 완성
\kfChoice 구성
\kfChoice 차례
\end{kfChoices}
\end{kfQBlock}
\begin{kfQBlock}{05}{서술형}
\kfQStem{'작성'을 사용한 문장}
\kfAnswerNote{답안}{4}
\end{kfQBlock}
\begin{kfQBlock}{06}{서술형}
\kfQStem{'구조'를 사용한 문장}
\kfAnswerNote{답안}{4}
\end{kfQBlock}


\kfStartGrammar{글의 짜임 (처음-가운데-끝)}


\kfPassageNoteNT{%
처음 (머리말): 글의 시작 부분. 읽는 사람의 관심을 끌고 무엇에 대해 쓸지 소개합니다.

가운데 (본문): 글의 중심 부분. 하고 싶은 말에 대한 구체적인 사실, 이유, 설명 등을 씁니다.

끝 (맺음말): 글의 마무리 부분. 전체 내용을 정리하고 자신의 생각이나 느낌을 씁니다.
\par\medskip\begin{itemize}[leftmargin=18pt, itemsep=2pt, topsep=2pt, label=\textbullet]
\item 처음(머리말) — 읽는 사람의 관심을 끌고 무엇에 대해 쓸지 소개
\item 가운데(본문) — 하고 싶은 말에 대한 구체적인 사실·이유·설명 제시
\item 끝(맺음말) — 전체 내용 정리와 자신의 생각·느낌 표현
\item 비유: 개요는 건물을 지을 때 필요한 설계도와 같다
\end{itemize}
}

\kfActivityBg \kfSecHeader{kfAreaGrammar}{활동}{지시어 연결}
\par\kfMaybeNeedspace{18\baselineskip}\noindent{\kfTipMark\,\,}{\small\color{kfInk} 지시어가 가리키는 대상을 찾아 답해 보세요.}\par\nopagebreak[4]\smallskip\nopagebreak[4]
\begin{enumerate}[leftmargin=22pt, itemsep=6pt, topsep=4pt, label=\textbf{\arabic*.}]
\item 다음 문장에서 '이것'과 '이러한 과정'이 가리키는 것을 찾아 쓰세요.
\begin{kfEvidence}"글을 쓰기 전에 쓸 내용을 미리 간추려 정리한 것을 개요라고 합니다. 개요는 글의 뼈대와 같습니다. 뼈대가 튼튼해야 우리 몸이 바로 서듯, \uline{이것}이 잘 짜여 있어야 좋은 글이 됩니다. \uline{이러한 과정}은 글쓰기 능력을 향상시키는 데 큰 도움이 됩니다."\end{kfEvidence}
\end{enumerate}

\kfActivityBg \par\kfMaybeNeedspace{24\baselineskip}\noindent{\kfTipMark\,\,}{\small\color{kfInk} [1-3] 다음 내용이 글의 어느 부분에 들어갈지 선으로 연결해 보세요.}\par\nopagebreak[4]\smallskip\nopagebreak[4]
\begin{kfConnect}
\kfPair{"지난 주말에 다녀온 놀이공원 이야기를 하려고 한다."}{㉠ 끝}
\kfPair{"정말 신나고 재미있는 하루였다. 다음에 또 가고 싶다."}{㉡ 가운데}
\kfPair{"놀이공원에서 롤러코스터를 탔는데 정말 짜릿했다."}{㉢ 처음}
\end{kfConnect}

\kfActivityBg \kfSecHeader{kfAreaGrammar}{활동}{어법 훈련}
\par\kfMaybeNeedspace{24\baselineskip}\noindent{\kfTipMark\,\,}{\small\color{kfInk} [1-2] 띄어쓰기를 잘못하면 뜻이 완전히 달라질 수 있습니다. 문장의 뜻이 잘 통하도록 띄어쓰기를 바르게 고쳐 써 보세요.}\par\nopagebreak[4]\smallskip\nopagebreak[4]
\begin{enumerate}[leftmargin=22pt, itemsep=6pt, topsep=4pt, label=\textbf{\arabic*.}]
\item 아버지가방에들어가신다. (아버지가 방에 들어가는 내용으로)
\item 나물좀다오. (마시는 물을 달라는 내용으로)
\end{enumerate}


\kfStartConcept{글쓰기에도 지도가 필요해!}


\kfPassageNoteNT{%
　낯선 곳으로 여행을 떠날 때 지도가 없으면 길을 잃기 쉽습니다. 글쓰기도 마찬가지입니다. 아무 계획 없이 글을 쓰기 시작하면, 내가 무슨 말을 하고 싶은지 헷갈리거나 내용이 엉뚱한 곳으로 흘러가기 쉽습니다. 이때 필요한 것이 바로 ‘개요’입니다.

　개요는 건물을 지을 때 필요한 설계도와 같습니다. 건축가가 설계도 없이 집을 지으면 기둥이 삐뚤어지거나 지붕이 무너질 수 있듯이, 글도 뼈대가 되는 개요가 없으면 내용이 뒤죽박죽 섞여버립니다. 글을 쓰기 전에 처음, 가운데, 끝에 어떤 내용을 넣을지 미리 짜임을 생각하고 순서를 정해보세요. 튼튼한 설계도가 멋진 집을 만들듯, 잘 짜인 개요는 훌륭한 글을 완성하게 해줍니다.
\par\medskip\begin{itemize}[leftmargin=18pt, itemsep=2pt, topsep=2pt, label=\textbullet]
\item 처음 — 글의 화제를 제시한다(예: '우리 학교 운동장')
\item 가운데 — 중심 내용을 단락별로 배치한다(시설·추억·문제점)
\item 끝 — 주장·감상으로 마무리한다('더 자주 가꾸어 쓰고 싶다')
\end{itemize}
}

\kfActivityBg \par\kfMaybeNeedspace{24\baselineskip}\noindent{\kfTipMark\,\,}{\small\color{kfInk} [1-3] 위 글의 내용과 일치하면 O, 일치하지 않으면 X 표시를 하세요.}\par\nopagebreak[4]\smallskip\nopagebreak[4]
\begin{kfOXList}
\kfOXItem 개요는 여행할 때 필요한 지도나 건물을 지을 때 필요한 설계도와 비슷한 역할을 한다. ( ○ /  ) \kfOX
\kfOXItem 글을 다 쓰고 나서 내용을 확인하기 위해 마지막에 개요를 작성한다. ( ○ /  ) \kfOX
\kfOXItem 개요를 미리 짜면 글의 내용이 엉뚱한 곳으로 빠지지 않도록 도와준다. ( ○ /  ) \kfOX
\end{kfOXList}


\kfStartLit{「민준이는 학교에서 돌아오자마자 블록 상자…」}


\kfPassageNote{}{%
　민준이는 학교에서 돌아오자마자 블록 상자를 바닥에 쏟았습니다.

"오늘은 세상에서 제일 높은 성을 쌓을 거야!"

　민준이는 마음이 급해서 눈에 보이는 블록을 닥치는 대로 쌓아 올리기 시작했습니다. 바닥은 좁고, 그 위에 큰 블록을 올리기도 했습니다.

　처음에는 그럴듯해 보였습니다. 하지만 성이 민준이의 무릎 높이만큼 올라갔을 때였습니다. 기둥이 휘청거리더니, '와르르!' 성은 순식간에 무너지고 말았습니다.

"아, 뭐야! 또 무너졌어."

　민준이가 속상해서 울상을 짓자, 옆에서 책을 읽던 형이 다가왔습니다.

"민준아, 설계도도 없이 무작정 쌓으니까 무너지는 거야. 집을 지을 때도 순서가 있단다."

　형은 종이와 연필을 가져와서 민준이 앞에 펼쳤습니다.

"자, 먼저 어떤 모양의 성을 만들지 그림을 그려보자. 바닥은 넓고 튼튼하게, 기둥은 반듯하게, 지붕은 뾰족하게 그리는 거야. 그게 바로 설계도야."

　민준이는 형과 함께 튼튼한 성의 모습을 그렸습니다. 그리고 그 그림을 보며 바닥부터 차근차근 쌓았습니다. 넓은 블록으로 기초를 다지고, 기둥을 세우고, 마지막으로 지붕을 조심스럽게 올렸습니다. 이번에는 아무리 높이 쌓아도 흔들리지 않았습니다.

"우와! 형, 이번에는 안 무너졌어!"

　민준이는 튼튼하게 완성된 성을 보며 활짝 웃었습니다.

　형이 민준이의 머리를 쓰다듬으며 말했습니다.

"글쓰기도 이것과 똑같아. 무작정 쓰기 시작하면 내용이 뒤죽박죽되고 무너지기 쉬워. 하지만 무엇을 쓸지 미리 계획을 세우면, 무너지지 않는 멋진 글을 쓸 수 있단다."
}


\kfActivityBg \par\kfMaybeNeedspace{24\baselineskip}\noindent{\kfTipMark\,\,}{\small\color{kfInk} [1-3] 이야기의 내용과 일치하면 O, 일치하지 않으면 X 표시를 하세요.}\par\nopagebreak[4]\smallskip\nopagebreak[4]
\begin{kfOXList}
\kfOXItem 민준이는 처음부터 형의 도움을 받아 성을 쌓았다. \kfOX
\kfOXItem 설계도를 그리고 쌓은 성은 높이 쌓아도 무너지지 않았다. \kfOX
\kfOXItem 형은 글쓰기에는 계획이 필요 없다고 말했다. \kfOX
\end{kfOXList}

\kfActivityBg \par\kfMaybeNeedspace{18\baselineskip}\noindent{\kfTipMark\,\,}{\small\color{kfInk} [1-3] 여러분이 '주말에 있었던 일'을 주제로 글을 쓴다면 어떤 설계도를 그릴지 상상해 보세요. 간단히 내용을 채워보세요.}\par\nopagebreak[4]\smallskip\nopagebreak[4]
\kfWritingPrompt{처음 (시작):}
\kfWriting{답안 작성}{8}
\kfWritingPrompt{가운데 (자세한 이야기):}
\kfWriting{답안 작성}{8}
\kfWritingPrompt{끝 (마무리):}
\kfWriting{답안 작성}{8}

\kfActivityBg \par\kfMaybeNeedspace{18\baselineskip}\noindent{\kfTipMark\,\,}{\small\color{kfInk} 다음 조건에 맞게 글을 써 보세요.}\par\nopagebreak[4]\smallskip\nopagebreak[4]
\kfWritingPrompt{계획(설계도, 개요)을 세워서 성공했던 경험이나, 계획 없이 해서 실패했던 경험이 있나요? 간단히 써 보세요.}
\kfWriting{답안 작성}{8}

\kfSecHeader{kfAreaLiterature}{문제}{}

\begin{kfQBlock}{01}{객관식}
\kfQStem{민준이가 처음에 쌓은 성이 무너진 이유는 무엇인가요?}
\begin{kfChoices}
\kfChoice 블록의 개수가 부족해서
\kfChoice 바람이 너무 세게 불어서
\kfChoice 계획 없이 급하게 쌓아서
\kfChoice 형이 실수로 건드려서
\end{kfChoices}
\end{kfQBlock}
\begin{kfQBlock}{02}{객관식}
\kfQStem{형이 민준이에게 성을 쌓기 전에 가장 먼저 하라고 한 것은 무엇인가요?}
\begin{kfChoices}
\kfChoice 블록을 색깔별로 분류하기
\kfChoice 어떤 모양을 만들지 그림 그리기
\kfChoice 더 크고 튼튼한 블록 사 오기
\kfChoice 친구들을 불러서 같이 쌓기
\end{kfChoices}
\end{kfQBlock}
\begin{kfQBlock}{03}{객관식}
\kfQStem{이 이야기에서 '블록 성 쌓기'는 무엇을 비유하고 있나요?}
\begin{kfChoices}
\kfChoice 그림 그리기
\kfChoice 글쓰기
\kfChoice 요리하기
\kfChoice 운동하기
\end{kfChoices}
\end{kfQBlock}
\begin{kfQBlock}{04}{객관식}
\kfQStem{성이 무너졌을 때와 완성되었을 때 민준이의 마음 변화로 알맞은 것은?}
\begin{kfChoices}
\kfChoice 기쁨 → 슬픔
\kfChoice 속상함 → 뿌듯함
\kfChoice 화남 → 걱정됨
\kfChoice 지루함 → 신남
\end{kfChoices}
\end{kfQBlock}
\begin{kfQBlock}{05}{객관식}
\kfQStem{형이 민준이에게 알려주고 싶었던 교훈은 무엇인가요?}
\begin{kfChoices}
\kfChoice 형제끼리는 사이좋게 지내야 한다.
\kfChoice 블록 놀이는 창의력을 키워준다.
\kfChoice 무엇이든 기초와 계획이 중요하다.
\kfChoice 실패는 성공의 어머니이다.
\end{kfChoices}
\end{kfQBlock}
\begin{kfQBlock}{06}{서술형}
\kfQStem{만약 민준이가 형의 조언을 듣지 않고 다시 그냥 쌓았다면 어떤 일이 벌어졌을지 상상하여 써보세요.}
\kfAnswerNote{답안}{4}
\end{kfQBlock}
\begin{kfQBlock}{07}{서술형}
\kfQStem{글쓰기에서 건물의 '설계도' 역할을 하는 것을 두 글자로 무엇이라고 할까요?}
\kfAnswerNote{답안}{4}
\end{kfQBlock}


\kfStartNonlit{개요를 작성하는 방법과 중요성}


\kfPassageNote{개요를 작성하는 방법과 중요성}{%
　글을 쓰기 전에 쓸 내용을 미리 간추려 정리한 것을 ‘개요’라고 합니다. 개요는 글의 뼈대와 같습니다. 뼈대가 튼튼해야 우리 몸이 바로 서듯, 개요가 잘 짜여 있어야 좋은 글이 됩니다. 개요를 작성하면 글의 흐름이 자연스러워지고, 꼭 써야 할 내용이 빠지거나 불필요한 내용이 중복되는 것을 막을 수 있습니다.

　개요를 작성할 때는 보통 ‘처음-가운데-끝’의 3단계 구성을 사용합니다.

　첫째, ‘처음’ 부분에서는 글을 쓰게 된 동기나 글에서 다룰 주제를 소개합니다. 읽는 사람의 흥미를 끌 수 있는 내용이면 더 좋습니다.

　둘째, ‘가운데’ 부분에서는 하고 싶은 말에 대한 구체적인 내용이나 근거, 예시를 들어 설명합니다. 글에서 가장 많은 부분을 차지하며, 내용을 알기 쉽게 차례대로 배치해야 합니다.

　셋째, ‘끝’ 부분에서는 전체 내용을 요약하거나 자신의 생각, 느낌을 정리하며 마무리합니다. 앞으로의 다짐이나 당부의 말을 쓰기도 합니다.

　개요를 짤 때는 먼저 떠오르는 생각들을 자유롭게 적어본 뒤, 관련된 내용끼리 분류하고, 글의 순서에 맞게 배열하는 과정을 거칩니다. 귀찮다고 생각하지 말고 5분만 투자하여 개요를 짜 보세요. 훨씬 더 쉽고 빠르게 멋진 글을 완성할 수 있을 것입니다.
}


\kfActivityBg \par\kfMaybeNeedspace{32\baselineskip}\noindent{\kfTipMark\,\,}{\small\color{kfInk} 다음은 각 문단의 핵심 내용을 정리한 것입니다. 빈칸에 알맞은 말을 채워 넣으세요.}\par\nopagebreak[4]\smallskip\nopagebreak[4]
\begin{kfStructure}
\kfNode{0}{}{}
\kfNode{0}{}{}
\kfNode{0}{}{}
\kfNode{0}{}{}
\end{kfStructure}

\kfActivityBg \par\kfMaybeNeedspace{18\baselineskip}\noindent{\kfTipMark\,\,}{\small\color{kfInk} [1-3] 여러분이 '가장 좋아하는 음식'을 소개하는 글을 쓴다고 생각하고 개요를 완성해 보세요.}\par\nopagebreak[4]\smallskip\nopagebreak[4]
\kfWritingPrompt{처음:}
\kfWriting{답안 작성}{8}
\kfWritingPrompt{가운데:}
\kfWriting{답안 작성}{8}
\kfWritingPrompt{끝:}
\kfWriting{답안 작성}{8}

\kfSecHeader{kfAreaNonlit}{문제}{}

\begin{kfQBlock}{01}{객관식}
\kfQStem{이 글의 주제로 가장 적절한 것은?}
\begin{kfChoices}
\kfChoice 글씨를 바르게 쓰는 방법
\kfChoice 개요의 뜻과 작성하는 방법
\kfChoice 재미있는 이야기를 지어내는 법
\kfChoice 친구와 사이좋게 지내는 법
\end{kfChoices}
\end{kfQBlock}
\begin{kfQBlock}{02}{객관식}
\kfQStem{개요를 작성했을 때의 장점이 \uline{\underline{아닌}} 것은?}
\begin{kfChoices}
\kfChoice 글의 흐름이 자연스러워진다.
\kfChoice 쓸 내용이 빠지는 것을 막을 수 있다.
\kfChoice 내용이 중복되는 것을 피할 수 있다.
\kfChoice 글의 분량이 무조건 길어진다.
\end{kfChoices}
\end{kfQBlock}
\begin{kfQBlock}{03}{객관식}
\kfQStem{글의 구성 단계 중 '가운데'에 들어갈 내용으로 알맞은 것은?}
\begin{kfChoices}
\kfChoice 글을 쓰게 된 동기
\kfChoice 전체 내용 요약
\kfChoice 구체적인 예시나 설명
\kfChoice 앞으로의 다짐
\end{kfChoices}
\end{kfQBlock}
\begin{kfQBlock}{04}{객관식}
\kfQStem{개요를 짤 때의 순서로 가장 자연스러운 것은?}
\begin{kfChoices}
\kfChoice 내용 배열하기 → 생각 떠올리기 → 관련 내용 묶기
\kfChoice 관련 내용 묶기 → 내용 배열하기 → 생각 떠올리기
\kfChoice 생각 떠올리기 → 관련 내용 묶기 → 내용 배열하기
\kfChoice 생각 떠올리기 → 내용 배열하기 → 관련 내용 묶기
\end{kfChoices}
\end{kfQBlock}
\begin{kfQBlock}{05}{서술형}
\kfQStem{글쓴이가 개요를 '뼈대'에 비유한 이유는 무엇일까요?}
\kfAnswerNote{답안}{4}
\end{kfQBlock}


\kfStartWeekly{실력확인 영역 학습}


\noindent\textit{\small (제한시간: 30분 / 총 15문항)}\par\medskip
\par\noindent\textbf{[4-6] 다음 글을 읽고 물음에 답하세요.}\par\medskip
\kfPassageNote{지문 4}{%
민준이는 처음에는 마음이 급해서 눈에 보이는 블록을 닥치는 대로 쌓아 올렸습니다. 하지만 성이 높아지자 기둥이 흔들리더니 순식간에 무너지고 말았습니다. 속상해하는 민준이에게 형이 다가와 말했습니다. "민준아, 집을 지을 때도 순서가 있단다. 먼저 어떤 모양으로 만들지 그림을 그려보자." 민준이는 형과 함께 튼튼한 성의 모습을 그렸습니다. 그리고 그 그림을 보며 바닥부터 차근차근 쌓았습니다. 넓은 블록으로 기초를 다지고, 기둥을 세우고, 마지막으로 지붕을 조심스럽게 올렸습니다. 이번에는 아무리 높이 쌓아도 흔들리지 않았습니다.
}
\par\noindent\textbf{[7-9] 다음 글을 읽고 물음에 답하세요.}\par\medskip
\kfPassageNote{지문 7}{%
글을 쓰기 전에 쓸 내용을 미리 간추려 정리한 것을 '개요'라고 합니다. 개요는 글의 뼈대와 같습니다. 뼈대가 튼튼해야 우리 몸이 바로 서듯, 개요가 잘 짜여 있어야 좋은 글이 됩니다. 개요를 작성할 때는 보통 '처음-가운데-끝'의 3단계 구성을 사용합니다. 첫째, '처음' 부분에서는 글을 쓰게 된 동기나 글에서 다룰 주제를 소개합니다. 둘째, '가운데' 부분에서는 하고 싶은 말에 대한 구체적인 내용이나 근거, 예시를 들어 설명합니다. 글에서 가장 많은 부분을 차지하며, 내용을 알기 쉽게 차례대로 배치해야 합니다. 셋째, '끝' 부분에서는 전체 내용을 요약하거나 자신의 생각, 느낌을 정리하며 마무리합니다.
}
\begin{kfQBlock}{01}{객관식}
\kfQStem{다음 중 '설계도'의 뜻으로 가장 알맞은 것은?}
\begin{kfChoices}
\kfChoice 글을 쓰기 전에 주요 내용을 간추려 짠 틀
\kfChoice 건축물의 구조나 계획을 그린 도면
\kfChoice 일정한 차례나 간격으로 늘어놓음
\kfChoice 어떤 물건의 몸체를 이루는 중요한 구조
\end{kfChoices}
\end{kfQBlock}

\begin{kfQBlock}{02}{객관식}
\kfQStem{다음 설명에 해당하는 낱말은 무엇인가요?}
\begin{kfEvidence}"부분들이 모여 전체를 이룬 짜임새"\end{kfEvidence}
\begin{kfChoices}
\kfChoice 배치
\kfChoice 분량
\kfChoice 구조
\kfChoice 차례
\end{kfChoices}
\end{kfQBlock}

\begin{kfQBlock}{03}{객관식}
\kfQStem{다음 빈칸에 들어갈 알맞은 말은?}
\begin{kfEvidence}"글을 쓰기 전에 개요를 [　　　　　]하면 훨씬 수월하게 글을 쓸 수 있다."\end{kfEvidence}
\begin{kfChoices}
\kfChoice 작성
\kfChoice 완성
\kfChoice 배치
\kfChoice 골조
\end{kfChoices}
\end{kfQBlock}

\begin{kfQBlock}{04}{객관식}
\kfQStem{위 글에서 민준이가 처음에 쌓은 성이 무너진 이유는 무엇인가요?}
\begin{kfChoices}
\kfChoice 블록의 색깔이 마음에 들지 않아서
\kfChoice 계획 없이 급하게 쌓았기 때문에
\kfChoice 형이 옆에서 방해를 했기 때문에
\kfChoice 바닥이 미끄러워서
\end{kfChoices}
\end{kfQBlock}

\begin{kfQBlock}{05}{객관식}
\kfQStem{민준이가 두 번째로 성을 쌓을 때 가장 먼저 한 일은 무엇인가요?}
\begin{kfChoices}
\kfChoice 지붕을 먼저 만들었다.
\kfChoice 더 많은 블록을 가져왔다.
\kfChoice 어떤 모양으로 만들지 그림을 그렸다.
\kfChoice 친구에게 도움을 청했다.
\end{kfChoices}
\end{kfQBlock}

\begin{kfQBlock}{06}{객관식}
\kfQStem{이 글의 중심 내용으로 가장 알맞은 것은?}
\begin{kfChoices}
\kfChoice 우리 몸의 뼈대가 하는 역할
\kfChoice 글의 종류와 특징
\kfChoice 개요의 뜻과 작성 방법
\kfChoice 재미있는 이야기를 짓는 법
\end{kfChoices}
\end{kfQBlock}

\begin{kfQBlock}{07}{객관식}
\kfQStem{위 글에서 설명하는 '가운데' 부분에 들어갈 내용으로 적절하지 \uline{\underline{않은}} 것은?}
\begin{kfChoices}
\kfChoice 하고 싶은 말에 대한 구체적인 내용
\kfChoice 주장을 뒷받침하는 근거
\kfChoice 이해를 돕기 위한 예시
\kfChoice 전체 내용을 요약하며 마무리하는 말
\end{kfChoices}
\end{kfQBlock}

\begin{kfQBlock}{08}{객관식}
\kfQStem{다음 문장에서 '무엇이(주어)'와 '어찌하다(서술어)'를 바르게 연결한 것은?}
\begin{kfEvidence}"튼튼한 설계도가 멋진 집을 만듭니다."\end{kfEvidence}
\begin{kfChoices}
\kfChoice 튼튼한 - 설계도가
\kfChoice 설계도가 - 만듭니다
\kfChoice 멋진 - 집을
\kfChoice 집을 - 만듭니다
\end{kfChoices}
\end{kfQBlock}

\begin{kfQBlock}{09}{객관식}
\kfQStem{글의 짜임 중 '처음' 부분의 역할로 알맞은 것은?}
\begin{kfChoices}
\kfChoice 구체적인 사실이나 이유를 든다.
\kfChoice 자신의 생각이나 느낌을 정리한다.
\kfChoice 읽는 사람의 흥미를 끌고 쓸 내용을 소개한다.
\kfChoice 전체 내용을 요약한다.
\end{kfChoices}
\end{kfQBlock}

\begin{kfQBlock}{10}{객관식}
\kfQStem{다음 중 글의 '끝' 부분에 들어가기에 가장 \uline{어색한} 내용은?}
\begin{kfChoices}
\kfChoice 앞으로의 다짐이나 계획
\kfChoice 전체 내용 요약
\kfChoice 글을 쓰게 된 동기나 계기
\kfChoice 글을 쓰고 난 후의 느낌
\end{kfChoices}
\end{kfQBlock}

\begin{kfQBlock}{11}{객관식}
\kfQStem{다음 중 띄어쓰기가 바르게 된 문장은?}
\begin{kfChoices}
\kfChoice 하늘이참 맑습니다.
\kfChoice 하늘 이 참 맑습니다.
\kfChoice 하늘이 참 맑습니다.
\kfChoice 하늘이 참맑습니다.
\end{kfChoices}
\end{kfQBlock}

\begin{kfQBlock}{12}{서술형}
\kfQStem{민준이가 완성한 성을 보고 형이 해 준 말의 핵심 내용을 한 문장으로 요약하여 쓰세요.}
\kfAnswerNote{답안}{4}
\end{kfQBlock}

\begin{kfQBlock}{13}{서술형}
\kfQStem{개요를 작성했을 때 얻을 수 있는 좋은 점을 한 가지 쓰세요.}
\kfAnswerNote{답안}{4}
\end{kfQBlock}

\begin{kfQBlock}{14}{서술형}
\kfQStem{다음 괄호 안에 들어갈 알맞은 말을 쓰세요.}
\begin{kfEvidence}"글을 쓸 때 [　　　　　]을/를 작성하면 글의 흐름이 자연스러워지고, 빠뜨리는 내용 없이 글을 쓸 수 있다."\end{kfEvidence}
\kfAnswerNote{답안}{4}
\end{kfQBlock}

\begin{kfQBlock}{15}{서술형}
\kfQStem{다음 문장의 뜻이 통하도록 띄어쓰기를 고쳐 쓰세요.}
\begin{kfEvidence}"오늘밤나무를심자." (밤에 나무를 심는다는 뜻으로)\end{kfEvidence}
\kfAnswerNote{답안}{4}
\end{kfQBlock}




\clearpage

\kfChapterCover{3}{프레게1 (챕터3)}{Chapter 3}{이번 챕터의 학습 내용을 시작합니다.}
\begin{kfChapterAreaList}
\kfChapterAreaItem{어휘}{kfAreaVocab}{}
\kfChapterAreaItem{문법}{kfAreaGrammar}{문장을 이어 주는 말 '접속사 (1)'}
\kfChapterAreaItem{개념}{kfAreaConcept}{긴 글을 한 입 크기로! '요약'의 힘}
\kfChapterAreaItem{문학}{kfAreaLiterature}{「"있잖아, 내가 지난주 토요일에 놀이공원…」}
\kfChapterAreaItem{비문학}{kfAreaNonlit}{문단별 중심 내용을 요약하는 기술}
\kfChapterAreaItem{실력확인}{kfAreaWeekly}{}
\end{kfChapterAreaList}

\kfStartVocab{어휘 영역 학습}


\kfActivityBg \par\kfMaybeNeedspace{24\baselineskip}\noindent{\kfTipMark\,\,}{\small\color{kfInk} [1-4] 다음 빈칸에 알맞은 어휘를 넣으세요.}\par\nopagebreak[4]\smallskip\nopagebreak[4]
\begin{kfBlankList}
\kfBlankItem 긴 이야기를 짧게 [　　　　　]하여 친구에게 들려주었다. (말이나 글의 요점을 잡아서 간추림을 나타내는 어휘를 넣으세요)
\kfBlankItem 글을 읽을 때는 가장 중요한 [　　　　　] 내용이 무엇인지 찾아야 한다. (사물의 가장 중심이 되는 중요한 부분을 나타내는 어휘를 넣으세요)
\kfBlankItem 이 글은 세 개의 [　　　　　]으로 이루어져 있다. (중심 생각 하나를 나타내는 문장들의 묶음을 나타내는 어휘를 넣으세요)
\kfBlankItem 그리고, 그러나, 그래서와 같은 말을 [　　　　　]라고 한다. (문장과 문장, 단어와 단어를 이어 주는 말을 나타내는 어휘를 넣으세요)
\end{kfBlankList}

\kfSecHeader{kfAreaVocab}{문제}{}

\begin{kfQBlock}{01}{객관식}
\kfQStem{'간추리다'와 가장 비슷한 의미를 가진 단어는?}
\begin{kfChoices}
\kfChoice 늘리다
\kfChoice 정리하다
\kfChoice 복사하다
\kfChoice 잊어버리다
\end{kfChoices}
\end{kfQBlock}
\begin{kfQBlock}{02}{객관식}
\kfQStem{'불필요'의 반의어로 가장 적절한 것은?}
\begin{kfChoices}
\kfChoice 중요
\kfChoice 필요
\kfChoice 소중
\kfChoice 부족
\end{kfChoices}
\end{kfQBlock}
\begin{kfQBlock}{03}{객관식}
\kfQStem{다음 문장에서 밑줄 친 부분과 바꿔 쓸 수 있는 말은?}
\begin{kfEvidence}"상대방의 의도를 정확하게 \uline{파악}하는 것이 중요하다."\end{kfEvidence}
\begin{kfChoices}
\kfChoice 이해
\kfChoice 오해
\kfChoice 무시
\kfChoice 질문
\end{kfChoices}
\end{kfQBlock}
\begin{kfQBlock}{04}{객관식}
\kfQStem{'중심'과 '세부'의 관계와 가장 비슷한 것은?}
\begin{kfChoices}
\kfChoice 핵심 - 요약
\kfChoice 전체 - 부분
\kfChoice 시작 - 출발
\kfChoice 문단 - 문장
\end{kfChoices}
\end{kfQBlock}
\begin{kfQBlock}{05}{서술형}
\kfQStem{'생략'을 사용한 문장}
\kfAnswerNote{답안}{4}
\end{kfQBlock}
\begin{kfQBlock}{06}{서술형}
\kfQStem{'불필요'를 사용한 문장}
\kfAnswerNote{답안}{4}
\end{kfQBlock}


\kfStartGrammar{문장을 이어 주는 말 '접속사 (1)'}


\kfPassageNoteNT{%
문장과 문장, 또는 단어와 단어를 자연스럽게 이어 주는 말을 접속사라고 합니다.

그리고: 앞의 내용과 뒤의 내용이 비슷하거나 이어질 때 씁니다. (나열, 순서)

그러나 (하지만): 앞의 내용과 뒤의 내용이 서로 반대될 때 씁니다. (대조)

그래서: 앞의 내용이 원인이 되어 뒤의 내용(결과)이 생길 때 씁니다. (인과)
\par\medskip\begin{itemize}[leftmargin=18pt, itemsep=2pt, topsep=2pt, label=\textbullet]
\item 그리고 — 나열·순서: '비가 왔다. 그리고 바람도 불었다.'
\item 그러나(하지만) — 대조: '열심히 공부했다. 그러나 점수가 낮았다.'
\item 그래서 — 인과: '늦잠을 잤다. 그래서 학교에 늦었다.'
\end{itemize}
}

\kfActivityBg \kfSecHeader{kfAreaGrammar}{활동}{지시어 연결}
\par\kfMaybeNeedspace{18\baselineskip}\noindent{\kfTipMark\,\,}{\small\color{kfInk} 지시어가 가리키는 대상을 찾아 답해 보세요.}\par\nopagebreak[4]\smallskip\nopagebreak[4]
\begin{enumerate}[leftmargin=22pt, itemsep=6pt, topsep=4pt, label=\textbf{\arabic*.}]
\item 다음 문장에서 '이 기술들'이 가리키는 것을 찾아 쓰세요.
\begin{kfEvidence}"가장 먼저 해야 할 일은 각 문단의 중심 문장을 찾는 것입니다. 두 번째로 할 일은 불필요한 내용을 생략하는 것입니다. 마지막으로 골라낸 중심 문장들을 자연스럽게 연결해야 합니다 ... \uline{이 기술들}을 익혀두면 아무리 긴 글이라도 두려워하지 않고 핵심 내용을 쏙쏙 뽑아낼 수 있습니다."\end{kfEvidence}
\end{enumerate}

\kfActivityBg \par\kfMaybeNeedspace{24\baselineskip}\noindent{\kfTipMark\,\,}{\small\color{kfInk} [1-4] 다음 빈칸에 알맞은 접속사를 넣어 문장을 완성하세요. (그리고, 그러나, 그래서)}\par\nopagebreak[4]\smallskip\nopagebreak[4]
\begin{kfBlankList}
\kfBlankItem 나는 사과를 좋아한다. \kfFillBlank[12mm]{} 배도 좋아한다.
\kfBlankItem 어제는 비가 많이 왔다. \kfFillBlank[12mm]{} 우리는 축구를 하지 못했다.
\kfBlankItem 열심히 공부했다. \kfFillBlank[12mm]{} 시험 성적은 좋지 않았다.
\kfBlankItem 민호는 일기를 썼다. \kfFillBlank[12mm]{} 잠자리에 들었다.
\end{kfBlankList}

\kfActivityBg \kfSecHeader{kfAreaGrammar}{활동}{어법 훈련}
\par\kfMaybeNeedspace{24\baselineskip}\noindent{\kfTipMark\,\,}{\small\color{kfInk} [1-2] 같은 뜻을 가진 말이 반복되면 문장이 어색하고 길어집니다. 불필요하게 반복된 부분을 찾아 한 번만 사용하여 자연스러운 문장으로 고쳐 보세요.}\par\nopagebreak[4]\smallskip\nopagebreak[4]
\begin{enumerate}[leftmargin=22pt, itemsep=6pt, topsep=4pt, label=\textbf{\arabic*.}]
\item 넓은 운동장에서 축구를 차고 놀았다. ('축구'는 '차다'는 뜻을 포함하고 있음)
\item 지나가는 행인에게 길을 물어보았다. ('행인'은 '지나가는 사람'이라는 뜻)
\end{enumerate}


\kfStartConcept{긴 글을 한 입 크기로! '요약'의 힘}


\kfPassageNoteNT{%
　우리는 매일 엄청난 양의 정보를 접합니다. 뉴스, 유튜브 영상, 교과서의 내용까지 모든 것을 다 기억할 수는 없습니다. 이때 필요한 것이 바로 ‘요약’입니다. 요약은 긴 글이나 말에서 불필요한 가지는 생략하고, 가장 중요한 핵심 뼈대만 남기는 것을 말합니다.

　요약을 잘하면 두 가지가 좋아집니다. 첫째, 내용을 더 잘 기억할 수 있습니다. 복잡한 내용을 내 머릿속 서랍에 들어가기 좋게 정리하는 과정이기 때문입니다. 둘째, 다른 사람에게 내 생각을 전달할 때 훨씬 명확해집니다. 횡설수설하지 않고, "이 이야기의 포인트는 이거야!"라고 말할 수 있게 되죠. 글을 읽을 때 각 문단의 중심 내용을 찾고 그것을 연결해 보세요. 그것이 요약의 시작입니다.
\par\medskip\begin{itemize}[leftmargin=18pt, itemsep=2pt, topsep=2pt, label=\textbullet]
\item 요약 방법: 각 문단 중심 문장 고르기 → 이어 붙이기
\item 효과 1 — 내용이 오래 기억된다.
\item 효과 2 — 다른 사람에게 쉽게 설명할 수 있다.
\end{itemize}
}

\kfActivityBg \par\kfMaybeNeedspace{24\baselineskip}\noindent{\kfTipMark\,\,}{\small\color{kfInk} [1-3] 위 글의 내용과 일치하면 O, 일치하지 않으면 X 표시를 하세요.}\par\nopagebreak[4]\smallskip\nopagebreak[4]
\begin{kfOXList}
\kfOXItem 요약은 긴 내용에서 불필요한 부분을 생략하고 중요한 핵심만 남기는 것이다. ( ○ /  ) \kfOX
\kfOXItem 내용을 요약해서 정리하면 복잡한 정보를 기억하는 데 도움이 된다. ( ○ /  ) \kfOX
\kfOXItem 요약을 잘하기 위해서는 글자 하나도 빠뜨리지 않고 그대로 외워야 한다. ( ○ /  ) \kfOX
\end{kfOXList}


\kfStartLit{「"있잖아, 내가 지난주 토요일에 놀이공원…」}


\kfPassageNote{}{%
"있잖아, 내가 지난주 토요일에 놀이공원에 갔는데, 아침 8시에 일어났어. 아니, 8시 30분이었나? 아무튼 일어나서 엄마가 해주신 토스트를 먹었는데 잼이 딸기잼이었어. 그리고 파란색 옷을 입고 나갔는데, 버스 정류장에 사람이 진짜 많은 거야. 500번 버스를 탔는데 자리가 없어서 서서 갔어. 가는 길에 차도 엄청 막혔고..."

　민호는 신이 나서 주말에 있었던 일을 이야기했습니다. 하지만 이야기를 듣던 수진이는 점점 표정이 어두워지더니 하품을 크게 했습니다.

"아함\textasciitilde{}. 그래서 민호야, 놀이공원에서 무슨 일이 있었는데?"

　민호는 머리를 긁적이며 말했습니다.

"어? 아, 그러니까... 놀이공원에 도착해서 롤러코스터를 탔는데, 내 바로 앞에 앉은 사람이 가발을 쓰고 있었나 봐. 롤러코스터가 슝\textasciitilde{} 하고 내려갈 때 그 사람 가발이 날아가서 내 무릎 위에 툭 떨어졌어! 진짜 웃기지?"

　그제야 수진이가 깔깔 웃으며 말했습니다.

"푸하하! 진짜 웃긴다. 민호야, 처음부터 '롤러코스터 타다가 앞사람 가발이 날아왔어!'라고 말했으면 더 재미있었을 거야. 아침에 토스트 먹은 거랑 버스 탄 이야기는 불필요한 내용이라서 듣다가 지루했어."

　민호는 얼굴이 조금 빨개졌습니다.

"그렇구나. 나는 자세하게 말하는 게 좋은 줄 알았어."

"자세한 것도 좋지만, 중요한 핵심만 쏙쏙 뽑아서 말하면 친구들이 네 이야기를 더 잘 들어줄 거야. 마치 맛있는 알맹이만 남기고 껍질은 까서 먹는 귤처럼 말이야!"

　민호는 고개를 끄덕였습니다.

"알겠어. 다음부터는 껍질은 빼고 알맹이만 이야기할게!"
}


\kfActivityBg \par\kfMaybeNeedspace{24\baselineskip}\noindent{\kfTipMark\,\,}{\small\color{kfInk} [1-3] 이야기의 내용과 일치하면 O, 일치하지 않으면 X 표시를 하세요.}\par\nopagebreak[4]\smallskip\nopagebreak[4]
\begin{kfOXList}
\kfOXItem 민호는 처음부터 가발 이야기를 했다. \kfOX
\kfOXItem 수진이는 민호의 이야기가 처음부터 끝까지 재미있었다. \kfOX
\kfOXItem 민호는 앞으로 중요한 내용 위주로 말하기로 다짐했다. \kfOX
\end{kfOXList}

\kfActivityBg \par\kfMaybeNeedspace{18\baselineskip}\noindent{\kfTipMark\,\,}{\small\color{kfInk} [1-2] 민호처럼 친구에게 해주고 싶은 재미있는 경험이 있나요? 처음에는 생각나는 대로 다 적어보고, 그중에서 불필요한 부분은 지우고 핵심만 남겨보세요.}\par\nopagebreak[4]\smallskip\nopagebreak[4]
\kfWritingPrompt{원래 하고 싶은 말 (길게):}
\kfWriting{답안 작성}{8}
\kfWritingPrompt{요약한 말 (핵심만):}
\kfWriting{답안 작성}{8}

\kfActivityBg \par\kfMaybeNeedspace{18\baselineskip}\noindent{\kfTipMark\,\,}{\small\color{kfInk} [1-1] 민호의 경험을 참고하여, 중요한 내용을 요약하여 말하거나 글을 쓸 때 얻을 수 있는 긍정적인 효과를 두 가지 이상 써 보세요.}\par\nopagebreak[4]\smallskip\nopagebreak[4]
\kfWritingPrompt{효과 작성:}
\kfWriting{답안 작성}{8}

\kfSecHeader{kfAreaLiterature}{문제}{}

\begin{kfQBlock}{01}{객관식}
\kfQStem{수진이가 민호의 이야기를 듣다가 하품을 한 이유는 무엇인가요?}
\begin{kfChoices}
\kfChoice 민호의 목소리가 너무 작아서 잘 안 들려서
\kfChoice 민호가 거짓말을 하고 있다고 생각해서
\kfChoice 민호가 중요한 내용 없이 너무 길게 말해서
\kfChoice 수진이가 어젯밤에 잠을 못 자서 피곤해서
\end{kfChoices}
\end{kfQBlock}
\begin{kfQBlock}{02}{객관식}
\kfQStem{민호가 주말에 겪은 일 중 가장 흥미롭고 중요한 사건은 무엇인가요?}
\begin{kfChoices}
\kfChoice 아침에 늦게 일어나서 토스트를 먹은 일
\kfChoice 버스 정류장에 사람이 많았던 일
\kfChoice 놀이공원 가는 길에 차가 막힌 일
\kfChoice 롤러코스터에서 앞사람의 가발이 날아온 일
\end{kfChoices}
\end{kfQBlock}
\begin{kfQBlock}{03}{객관식}
\kfQStem{수진이는 '핵심만 말하는 것'을 무엇에 비유했나요?}
\begin{kfChoices}
\kfChoice 껍질을 까고 먹는 귤
\kfChoice 가지를 쳐낸 나무
\kfChoice 뼈대만 남은 건물
\kfChoice 엑기스만 뽑은 주스
\end{kfChoices}
\end{kfQBlock}
\begin{kfQBlock}{04}{객관식}
\kfQStem{다음 중 민호의 이야기에서 빼도 되는(불필요한) 내용은 무엇인가요? (2개 고르세요)}
\begin{kfChoices}
\kfChoice 아침에 늦게 일어나 토스트를 먹은 것
\kfChoice 놀이공원에서 롤러코스터를 탄 것
\kfChoice 버스 정류장에 사람이 많았던 것
\kfChoice 롤러코스터에서 앞사람의 가발이 날아온 것
\end{kfChoices}
\end{kfQBlock}
\begin{kfQBlock}{05}{객관식}
\kfQStem{친구에게 말을 하거나 글을 쓸 때 요약이 필요한 이유는 무엇일까요?}
\begin{kfChoices}
\kfChoice 상대방이 내 말을 지루해하지 않고 잘 이해할 수 있게 하기 위해서
\kfChoice 상대방에게 내가 아는 것이 많다는 것을 자랑하기 위해서
\kfChoice 이야기를 최대한 길게 늘려서 시간을 때우기 위해서
\kfChoice 중요한 비밀을 숨기기 위해서
\end{kfChoices}
\end{kfQBlock}
\begin{kfQBlock}{06}{서술형}
\kfQStem{민호가 처음에 했던 긴 이야기를 '한 문장'으로 요약해 보세요. (힌트: 누가 + 어디서 + 무엇을 + 어떻게 했다)}
\kfAnswerNote{답안}{4}
\end{kfQBlock}


\kfStartNonlit{문단별 중심 내용을 요약하는 기술}


\kfPassageNote{문단별 중심 내용을 요약하는 기술}{%
　글을 읽고 내용을 요약하는 것은 마치 여행 가방을 싸는 것과 비슷합니다. 여행을 갈 때 집안의 모든 물건을 가져갈 수 없듯이, 요약을 할 때도 글의 모든 내용을 다 쓸 수는 없습니다. 꼭 필요한 옷과 세면도구만 챙기듯, 글에서도 가장 중요한 핵심 내용만 골라내야 합니다. 그렇다면 긴 글을 어떻게 꼼꼼하고 간결하게 요약할 수 있을까요?

　가장 먼저 해야 할 일은 각 문단의 중심 문장을 찾는 것입니다. 하나의 문단은 보통 하나의 중심 생각과 그것을 설명하는 여러 뒷받침 문장으로 이루어져 있습니다. 중심 문장은 주로 문단의 맨 앞이나 맨 뒤에 있는 경우가 많습니다. 만약 중심 문장이 명확하지 않다면, 문단에서 가장 많이 반복되는 단어(핵심어)를 찾아 스스로 문장을 만들어 보아야 합니다.

　두 번째로 할 일은 불필요한 내용을 생략하는 것입니다. 구체적인 예시, 꾸며주는 말, 반복되는 설명은 과감하게 지워버리세요. 예를 들어, "사과, 배, 포도, 수박 같은 과일들"이라는 문장은 그냥 "여러 가지 과일"이라고 줄여서 표현할 수 있습니다.

　마지막으로 골라낸 중심 문장들을 자연스럽게 연결해야 합니다. 이때 '그리고', '그러나', '그래서'와 같은 접속사를 적절히 사용하면 문장들이 매끄럽게 이어져 하나의 완성된 요약글이 됩니다.

　이 기술들을 익혀두면 아무리 긴 글이라도 두려워하지 않고 핵심 내용을 쏙쏙 뽑아낼 수 있습니다. 요약하는 습관은 여러분의 독해력을 쑥쑥 키워주는 최고의 비결입니다.
}


\kfActivityBg \par\kfMaybeNeedspace{32\baselineskip}\noindent{\kfTipMark\,\,}{\small\color{kfInk} 다음은 각 문단의 핵심 내용을 정리한 것입니다. 빈칸에 알맞은 말을 채워 넣으세요.}\par\nopagebreak[4]\smallskip\nopagebreak[4]
\begin{kfStructure}
\kfNode{0}{}{}
\kfNode{0}{}{}
\kfNode{0}{}{}
\kfNode{0}{}{}
\kfNode{0}{}{}
\end{kfStructure}

\kfActivityBg \par\kfMaybeNeedspace{18\baselineskip}\noindent{\kfTipMark\,\,}{\small\color{kfInk} [1-1] 다음 문단을 읽고 불필요한 부분은 줄을 그어 지우고, 핵심 내용만 남겨 한 문장으로 요약해 보세요. "나는 어제 도서관에 가서 책을 빌렸다. 도서관에는 재미있는 만화책도 있었고, 공룡에 대한 책도 있었고, 위인전도 있었다. 나는 그중에서 평소에 읽고 싶었던 이순신 장군에 대한 위인전을 한 권 빌려서 집으로 돌아왔다."}\par\nopagebreak[4]\smallskip\nopagebreak[4]
\kfWritingPrompt{요약한 문장:}
\kfWriting{답안 작성}{8}

\kfSecHeader{kfAreaNonlit}{문제}{}

\begin{kfQBlock}{01}{객관식}
\kfQStem{이 글의 주제로 가장 적절한 것은?}
\begin{kfChoices}
\kfChoice 여행 가방을 효율적으로 싸는 법
\kfChoice 접속사를 사용하여 문장을 잇는 법
\kfChoice 문단별 중심 내용을 요약하는 방법
\kfChoice 글을 길고 자세하게 쓰는 방법
\end{kfChoices}
\end{kfQBlock}
\begin{kfQBlock}{02}{객관식}
\kfQStem{글쓴이가 제시한 요약의 방법이 \uline{\underline{아닌}} 것은?}
\begin{kfChoices}
\kfChoice 각 문단의 중심 문장을 정확히 찾는다.
\kfChoice 불필요한 예시나 자세한 설명은 생략한다.
\kfChoice 접속사를 사용하여 문장들을 연결한다.
\kfChoice 모든 문장을 빠짐없이 그대로 옮겨 적는다.
\end{kfChoices}
\end{kfQBlock}
\begin{kfQBlock}{03}{객관식}
\kfQStem{다음 <보기>의 문단을 읽고 중심 내용을 가장 잘 요약한 것을 고르세요.}
\begin{kfEvidence}"여름철에는 음식이 쉽게 상할 수 있습니다. 높은 온도와 습도 때문에 세균이 빨리 번식하기 때문입니다. 따라서 남은 음식은 반드시 냉장고에 보관해야 하고, 유통기한을 꼼꼼히 확인하는 것이 좋습니다."\end{kfEvidence}
\begin{kfChoices}
\kfChoice 여름철 날씨는 덥고 습하다.
\kfChoice 세균은 높은 온도에서 빨리 번식한다.
\kfChoice 유통기한이 지난 음식은 먹으면 안 된다.
\kfChoice 여름철에는 음식이 상하기 쉬우므로 보관에 주의해야 한다.
\end{kfChoices}
\end{kfQBlock}
\begin{kfQBlock}{04}{객관식}
\kfQStem{요약하는 과정에서 '불필요한 내용 생략'에 대한 설명으로 알맞지 \uline{\underline{않은}} 것은?}
\begin{kfChoices}
\kfChoice 구체적인 예시와 설명은 줄여야 한다.
\kfChoice 반복되는 내용은 오히려 강조해야 한다.
\kfChoice 꾸며주는 말은 요약 시 지우는 것이 좋다.
\kfChoice 글의 뼈대만 남기기 위해 과감히 지운다.
\end{kfChoices}
\end{kfQBlock}
\begin{kfQBlock}{05}{서술형}
\kfQStem{글쓴이가 요약할 때 '접속사'를 사용하라고 한 이유는 무엇인가요?}
\kfAnswerNote{답안}{4}
\end{kfQBlock}


\kfStartWeekly{실력확인 영역 학습}


\noindent\textit{\small (제한시간: 30분 / 총 15문항)}\par\medskip
\par\noindent\textbf{[4-6] 다음 글을 읽고 물음에 답하세요.}\par\medskip
\kfPassageNote{지문 4}{%
민호는 신이 나서 주말에 있었던 일을 이야기했습니다. 하지만 이야기를 듣던 수진이는 점점 표정이 어두워지더니 하품을 크게 했습니다. "아함\textasciitilde{}. 그래서 민호야, 놀이공원에서 무슨 일이 있었는데?" 민호는 머리를 긁적이며 말했습니다. "어? 아, 그러니까... 롤러코스터를 탔는데, 내 바로 앞에 앉은 사람이 가발을 쓰고 있었나 봐. 롤러코스터가 슝\textasciitilde{} 하고 내려갈 때 그 사람 가발이 날아가서 내 무릎 위에 툭 떨어졌어! 진짜 웃기지?" 그제야 수진이가 깔깔 웃으며 말했습니다. "푸하하! 진짜 웃긴다. 민호야, 처음부터 '롤러코스터 타다가 앞사람 가발이 날아왔어!'라고 말했으면 더 재미있었을 거야. ( ㉠ )한 내용이라서 듣다가 지루했어."
}
\par\noindent\textbf{[7-9] 다음 글을 읽고 물음에 답하세요.}\par\medskip
\kfPassageNote{지문 7}{%
글을 읽고 내용을 요약하는 것은 마치 여행 가방을 싸는 것과 비슷합니다. 꼭 필요한 물건만 챙기듯, 글에서도 가장 중요한 핵심 내용만 골라내야 합니다. 요약을 잘하려면 첫째, 각 문단의 중심 문장을 찾아야 합니다. 중심 문장은 주로 문단의 맨 앞이나 맨 뒤에 있는 경우가 많습니다. 둘째, 불필요한 내용은 과감하게 ( ㉡ )해야 합니다. 구체적인 예시나 꾸며주는 말, 반복되는 설명은 지워버리세요. 셋째, 골라낸 중심 문장들을 접속사를 사용하여 자연스럽게 연결해야 합니다.
}
\begin{kfQBlock}{01}{객관식}
\kfQStem{다음 중 '요약'의 뜻으로 가장 알맞은 것은?}
\begin{kfChoices}
\kfChoice 글자 수를 세어서 기록하는 것
\kfChoice 글의 내용을 빠짐없이 그대로 베껴 쓰는 것
\kfChoice 말이나 글의 요점을 잡아 간추리는 것
\kfChoice 자신의 상상력을 더해 새로운 이야기를 만드는 것
\end{kfChoices}
\end{kfQBlock}

\begin{kfQBlock}{02}{객관식}
\kfQStem{다음 중 낱말의 관계가 나머지와 \uline{다른} 하나는?}
\begin{kfChoices}
\kfChoice 필요 - 불필요
\kfChoice 중심 - 세부
\kfChoice 생략 - 삭제
\kfChoice 시작 - 끝
\end{kfChoices}
\end{kfQBlock}

\begin{kfQBlock}{03}{객관식}
\kfQStem{위 글에서 수진이가 민호의 이야기를 듣다가 하품을 한 이유는 무엇인가요?}
\begin{kfChoices}
\kfChoice 민호의 목소리가 너무 작아서 잘 들리지 않아서
\kfChoice 민호가 거짓말을 하고 있다고 생각해서
\kfChoice 중요한 내용 없이 쓸데없는 말을 너무 길게 해서
\kfChoice 수진이가 어젯밤에 잠을 못 자서 피곤해서
\end{kfChoices}
\end{kfQBlock}

\begin{kfQBlock}{04}{객관식}
\kfQStem{위 글의 빈칸 ( ㉠ )에 들어갈 알맞은 말은 무엇인가요?}
\begin{kfChoices}
\kfChoice 핵심
\kfChoice 불필요
\kfChoice 감동적
\kfChoice 슬픈
\end{kfChoices}
\end{kfQBlock}

\begin{kfQBlock}{05}{객관식}
\kfQStem{위 글에서 글쓴이가 요약하는 방법을 설명하기 위해 빗댄 것은 무엇인가요?}
\begin{kfChoices}
\kfChoice 요리하기
\kfChoice 여행 가방 싸기
\kfChoice 그림 그리기
\kfChoice 청소하기
\end{kfChoices}
\end{kfQBlock}

\begin{kfQBlock}{06}{객관식}
\kfQStem{위 글에 따르면 요약을 할 때 접속사를 사용하는 이유는 무엇인가요?}
\begin{kfChoices}
\kfChoice 문장을 더 길게 만들기 위해서
\kfChoice 중심 문장들을 자연스럽게 연결하기 위해서
\kfChoice 어려운 단어를 쉽게 풀이하기 위해서
\kfChoice 글씨를 예쁘게 쓰기 위해서
\end{kfChoices}
\end{kfQBlock}

\begin{kfQBlock}{07}{객관식}
\kfQStem{다음 중 접속사에 대한 설명으로 알맞지 \uline{\underline{않은}} 것은?}
\begin{kfChoices}
\kfChoice 문장과 문장을 이어 주는 역할을 한다.
\kfChoice '그리고'는 앞뒤 내용이 비슷하거나 이어질 때 쓴다.
\kfChoice '그러나'는 앞뒤 내용이 서로 반대될 때 쓴다.
\kfChoice '그래서'는 앞뒤 내용이 아무런 관계가 없을 때 쓴다.
\end{kfChoices}
\end{kfQBlock}

\begin{kfQBlock}{08}{객관식}
\kfQStem{다음 빈칸에 공통으로 들어갈 알맞은 접속사는?}
\begin{kfEvidence}비가 많이 왔다. [　　　　　] 소풍이 취소되었다. \\\relax
늦잠을 잤다. [　　　　　] 지각을 하고 말았다.\end{kfEvidence}
\begin{kfChoices}
\kfChoice 하지만
\kfChoice 그러나
\kfChoice 그래서
\kfChoice 또는
\end{kfChoices}
\end{kfQBlock}

\begin{kfQBlock}{09}{객관식}
\kfQStem{다음 중 의미가 중복되지 않고 간결하게 잘 쓴 문장은?}
\begin{kfChoices}
\kfChoice 나는 넓은 운동장에서 축구를 찼다.
\kfChoice 지나가는 행인에게 길을 물어보았다.
\kfChoice 어제 도서관에서 책을 빌렸다.
\kfChoice 따뜻한 온수를 마셨다.
\end{kfChoices}
\end{kfQBlock}

\begin{kfQBlock}{10}{서술형}
\kfQStem{다음 설명에 해당하는 낱말을 쓰세요.}
\begin{kfEvidence}"사물의 가장 중심이 되는 중요한 부분"\end{kfEvidence}
\kfAnswerNote{답안}{4}
\end{kfQBlock}

\begin{kfQBlock}{11}{서술형}
\kfQStem{민호가 수진이에게 전하고 싶었던 이야기의 핵심은 무엇인지 한 문장으로 쓰세요.}
\kfAnswerNote{답안}{4}
\end{kfQBlock}

\begin{kfQBlock}{12}{서술형}
\kfQStem{위 글의 빈칸 ( ㉡ )에 들어갈 알맞은 낱말을 쓰세요.}
\kfAnswerNote{답안}{4}
\end{kfQBlock}

\begin{kfQBlock}{13}{서술형}
\kfQStem{다음 문장에서 '누가(주어)'와 '어찌하다(서술어)'를 찾아 쓰세요.}
\begin{kfEvidence}"민호는 수진이에게 재미있는 이야기를 해주었습니다."\end{kfEvidence}
\kfAnswerNote{답안}{4}
\end{kfQBlock}

\begin{kfQBlock}{14}{서술형}
\kfQStem{다음 문장의 괄호 안에 알맞은 접속사를 쓰세요.}
\begin{kfEvidence}"열심히 공부했다. [　　　　　] 시험 성적은 좋지 않았다."\end{kfEvidence}
\kfAnswerNote{답안}{4}
\end{kfQBlock}

\begin{kfQBlock}{15}{서술형}
\kfQStem{다음 문장에서 불필요하게 반복된 부분을 찾아 자연스럽게 고쳐 쓰세요.}
\begin{kfEvidence}"역전 앞에는 사람들이 많이 모여 있었다." ('역전'은 '역 앞'이라는 뜻)\end{kfEvidence}
\kfAnswerNote{답안}{4}
\end{kfQBlock}




\clearpage

\kfChapterCover{4}{프레게1 (챕터4)}{Chapter 4}{이번 챕터의 학습 내용을 시작합니다.}
\begin{kfChapterAreaList}
\kfChapterAreaItem{어휘}{kfAreaVocab}{}
\kfChapterAreaItem{문법}{kfAreaGrammar}{문장을 이어 주는 말 '접속사 (2)'}
\kfChapterAreaItem{개념}{kfAreaConcept}{내일 날씨는 어떻게 알 수 있을까?}
\kfChapterAreaItem{문학}{kfAreaLiterature}{「햇볕이 쨍쨍 내리쬐는 무더운 여름날이었습…」}
\kfChapterAreaItem{비문학}{kfAreaNonlit}{일기 예보는 어떻게 만들어질까요?}
\kfChapterAreaItem{실력확인}{kfAreaWeekly}{}
\end{kfChapterAreaList}

\kfStartVocab{어휘 영역 학습}


\kfActivityBg \par\kfMaybeNeedspace{24\baselineskip}\noindent{\kfTipMark\,\,}{\small\color{kfInk} [1-4] 다음 빈칸에 알맞은 어휘를 넣으세요.}\par\nopagebreak[4]\smallskip\nopagebreak[4]
\begin{kfBlankList}
\kfBlankItem 내일 비가 올 [　　　　　]은/는 80\%입니다. (어떤 일이 일어날 가능성을 나타내는 어휘를 넣으세요)
\kfBlankItem 태풍이 온다는 소식을 듣고 피해가 없도록 철저히 [　　　　　]했다. (앞으로 일어날 일에 대응하기 위해 미리 준비함을 나타내는 어휘를 넣으세요)
\kfBlankItem 인공위성으로 구름의 이동을 [　　　　　]한다. (자연 현상 등을 자세히 살펴봄을 나타내는 어휘를 넣으세요)
\kfBlankItem 기상청에서 이번 주말에 많은 눈이 내릴 것이라고 [　　　　　]했다. (앞으로 일어날 일을 미리 알림을 나타내는 어휘를 넣으세요)
\end{kfBlankList}

\kfSecHeader{kfAreaVocab}{문제}{}

\begin{kfQBlock}{01}{객관식}
\kfQStem{'예측'과 가장 비슷한 의미를 가진 단어는?}
\begin{kfChoices}
\kfChoice 예상
\kfChoice 결과
\kfChoice 후회
\kfChoice 확인
\end{kfChoices}
\end{kfQBlock}
\begin{kfQBlock}{02}{객관식}
\kfQStem{'대비'와 가장 비슷한 의미를 가진 단어는?}
\begin{kfChoices}
\kfChoice 걱정
\kfChoice 준비
\kfChoice 포기
\kfChoice 시작
\end{kfChoices}
\end{kfQBlock}
\begin{kfQBlock}{03}{객관식}
\kfQStem{다음 문장에서 밑줄 친 부분과 바꿔 쓸 수 있는 말은?}
\begin{kfEvidence}"과학자들은 수집한 정보를 자세히 \uline{분석}했다."\end{kfEvidence}
\begin{kfChoices}
\kfChoice 쪼개어 살폈다
\kfChoice 대충 보았다
\kfChoice 모아서 버렸다
\kfChoice 숨겼다
\end{kfChoices}
\end{kfQBlock}
\begin{kfQBlock}{04}{객관식}
\kfQStem{'습도'와 '축축함'의 관계와 가장 비슷한 것은?}
\begin{kfChoices}
\kfChoice 온도 - 따뜻함
\kfChoice 바람 - 멈춤
\kfChoice 구름 - 맑음
\kfChoice 태양 - 밤
\end{kfChoices}
\end{kfQBlock}
\begin{kfQBlock}{05}{서술형}
\kfQStem{'영향'을 사용한 문장}
\kfAnswerNote{답안}{4}
\end{kfQBlock}
\begin{kfQBlock}{06}{서술형}
\kfQStem{'현상'을 사용한 문장}
\kfAnswerNote{답안}{4}
\end{kfQBlock}


\kfStartGrammar{문장을 이어 주는 말 '접속사 (2)'}


\kfPassageNoteNT{%
왜냐하면: 앞의 내용에 대한 까닭(이유)을 말할 때 씁니다. 뒤에 '\textasciitilde{}때문이다'가 옵니다.

따라서: 앞의 내용이 원인이 되어 뒤에 결과나 결론이 나올 때 씁니다.

또는 (혹은): 여럿 중에서 하나를 선택할 때 씁니다.

반면에: 앞의 내용과 뒤의 내용이 서로 반대되거나 비교될 때 씁니다.
\par\medskip\begin{itemize}[leftmargin=18pt, itemsep=2pt, topsep=2pt, label=\textbullet]
\item 왜냐하면 — 이유: '나는 우산을 챙겼다. 왜냐하면 비가 올 것 같았기 때문이다.'
\item 따라서 — 결론: '연습을 꾸준히 했다. 따라서 좋은 결과를 얻었다.'
\item 또는(혹은) — 선택: '주스 또는 우유를 마실 수 있다.'
\item 반면에 — 대조: '형은 활발하다. 반면에 동생은 조용하다.'
\end{itemize}
}

\kfActivityBg \kfSecHeader{kfAreaGrammar}{활동}{지시어 연결}
\par\kfMaybeNeedspace{18\baselineskip}\noindent{\kfTipMark\,\,}{\small\color{kfInk} 지시어가 가리키는 대상을 찾아 답해 보세요.}\par\nopagebreak[4]\smallskip\nopagebreak[4]
\begin{enumerate}[leftmargin=22pt, itemsep=6pt, topsep=4pt, label=\textbf{\arabic*.}]
\item 다음 문장에서 '이것'이 가리키는 것을 찾아 쓰세요.
\begin{kfEvidence}"두 번째 단계는 '분석'입니다. 관측을 통해 모은 엄청난 양의 정보를 '슈퍼컴퓨터'에 입력합니다. 슈퍼컴퓨터는 복잡한 수학 공식으로 지구의 공기 흐름을 계산하여 미래의 날씨 변화를 그려냅니다. \uline{이것}을 '수치 예보 모델'이라고 합니다."\end{kfEvidence}
\end{enumerate}

\kfActivityBg \par\kfMaybeNeedspace{24\baselineskip}\noindent{\kfTipMark\,\,}{\small\color{kfInk} [1-4] 다음 빈칸에 알맞은 접속사를 넣어 문장을 완성하세요. (왜냐하면, 따라서, 또는, 반면에)}\par\nopagebreak[4]\smallskip\nopagebreak[4]
\begin{kfBlankList}
\kfBlankItem 여름에는 덥고 비가 많이 온다. \kfFillBlank[12mm]{} 겨울에는 춥고 건조하다.
\kfBlankItem 영수는 이번 추석 때 기차 \kfFillBlank[12mm]{} 고속버스로 시골에 내려갈 예정이다.
\kfBlankItem 기상청에서 태풍 경보를 내렸다. \kfFillBlank[12mm]{} 배들은 항구로 돌아와야 했다.
\kfBlankItem 개미는 쉴 수가 없었다. \kfFillBlank[12mm]{} 곧 겨울이 닥칠 것을 알았기 때문이다.
\end{kfBlankList}

\kfActivityBg \kfSecHeader{kfAreaGrammar}{활동}{어법 훈련}
\par\kfMaybeNeedspace{24\baselineskip}\noindent{\kfTipMark\,\,}{\small\color{kfInk} [1-3] 문장의 뜻에 알맞은 낱말을 골라 ○표 하세요.}\par\nopagebreak[4]\smallskip\nopagebreak[4]
\begin{enumerate}[leftmargin=22pt, itemsep=6pt, topsep=4pt, label=\textbf{\arabic*.}]
\item 기상청이 내일 날씨를 정확하게 ( 맞췄다 / 맞혔다 ).
\item 먹구름이 끼더니 ( 금세 / 금새 ) 비가 쏟아졌습니다.
\item 강한 바람에 창문이 쾅 ( 닫혔다 / 다쳤다 ).
\end{enumerate}


\kfStartConcept{내일 날씨는 어떻게 알 수 있을까?}


\kfPassageNoteNT{%
　옛날 사람들은 제비가 낮게 날면 비가 온다고도 생각했어요. 비가 오기 전에는 공기 상태가 달라져 곤충이 낮게 날 때가 있고, 제비도 먹이를 따라 낮게 날 수 있기 때문입니다. 하지만 오늘날에는 과학 기술의 발달로 훨씬 정확하게 날씨를 알 수 있습니다.

　우선 인공위성과 레이더가 하늘의 구름 이동, 바람의 세기, 대기의 상태 등을 실시간으로 관측합니다. 이렇게 수집된 엄청난 양의 정보는 슈퍼컴퓨터로 보내집니다. 슈퍼컴퓨터는 복잡한 계산을 통해 자료를 분석하고, 미래의 날씨 변화를 계산해 냅니다. 기상청 사람들은 이 결과를 바탕으로 내일 비가 올 확률이 얼마나 되는지, 기온은 몇 도일지 알려줍니다. 이것을 '일기 예보'라고 합니다. 일기 예보 덕분에 우리는 우산을 챙기거나 태풍에 미리 대비할 수 있습니다.
\par\medskip\begin{itemize}[leftmargin=18pt, itemsep=2pt, topsep=2pt, label=\textbullet]
\item 관측 도구 — 인공위성, 레이더
\item 분석 도구 — 슈퍼컴퓨터
\item 활용 — 우산 준비, 태풍 대비
\end{itemize}
}

\kfActivityBg \par\kfMaybeNeedspace{24\baselineskip}\noindent{\kfTipMark\,\,}{\small\color{kfInk} [1-3] 위 글의 내용과 일치하면 O, 일치하지 않으면 X 표시를 하세요.}\par\nopagebreak[4]\smallskip\nopagebreak[4]
\begin{kfOXList}
\kfOXItem 오늘날에는 슈퍼컴퓨터를 이용하여 옛날보다 훨씬 정확하게 날씨를 예측할 수 있다. ( ○ /  ) \kfOX
\kfOXItem 슈퍼컴퓨터는 인공위성과 레이더가 관측한 정보를 바탕으로 날씨 변화를 분석한다. ( ○ /  ) \kfOX
\kfOXItem 일기 예보는 단지 제비가 낮게 나는 것을 보고 비가 올 확률을 알려주는 것이다. ( ○ /  ) \kfOX
\end{kfOXList}


\kfStartLit{「햇볕이 쨍쨍 내리쬐는 무더운 여름날이었습…」}


\kfPassageNote{}{%
　햇볕이 쨍쨍 내리쬐는 무더운 여름날이었습니다. 베짱이는 나무 그늘 아래서 시원한 바람을 맞으며 바이올린을 켜고 노래를 불렀습니다.

"랄랄라\textasciitilde{} 날씨도 좋고 기분도 좋구나. 이렇게 좋은 날에는 신나게 놀아야지!"

　그때, 개미들이 줄을 지어 땀을 뻘뻘 흘리며 지나갔습니다. 개미들은 자기 몸집보다 훨씬 큰 먹이를 짊어지고 부지런히 집으로 나르고 있었습니다.

　베짱이가 개미를 보고 비웃으며 말했습니다.

"이봐, 개미 친구들! 이렇게 화창한 날에 왜 그렇게 힘들게 일을 해? 나랑 같이 노래나 부르며 놀자."

　맨 앞에 가던 개미가 잠시 멈춰 서서 땀을 닦으며 대답했습니다.

"우리는 다가올 비와 겨울을 대비해서 식량을 모으는 중이야. 하늘을 보니 곧 큰비가 쏟아질 것 같거든. 서쪽 하늘에 먹구름이 몰려오는 게 보여."

"하하하! 무슨 소리야? 하늘이 이렇게 맑고 푸른데 비가 온다고? 너희는 너무 걱정이 많아."

　베짱이는 개미의 말을 무시하고 다시 신나게 노래를 불렀습니다.

　하지만 며칠 뒤, 개미의 예측대로 하늘에 먹구름이 잔뜩 끼더니 굵은 빗방울이 쏟아지기 시작했습니다. 비는 며칠 동안 계속 내렸습니다. 성실하게 일하며 먹이를 모아둔 개미들은 집 안에서 배불리 먹으며 편안하게 지냈습니다.

　반면에 베짱이는 비를 피할 곳도 마땅치 않았고, 먹을 것이 없어 배가 고팠습니다. 덜덜 떨던 베짱이는 결국 개미네 집 문을 두드렸습니다.

"개미야, 나 좀 도와줘. 비가 이렇게 많이 올 줄 몰랐어."
}


\kfActivityBg \par\kfMaybeNeedspace{24\baselineskip}\noindent{\kfTipMark\,\,}{\small\color{kfInk} [1-3] 이야기의 내용과 일치하면 O, 일치하지 않으면 X 표시를 하세요.}\par\nopagebreak[4]\smallskip\nopagebreak[4]
\begin{kfOXList}
\kfOXItem 베짱이는 개미의 말을 듣고 함께 일을 도왔다. \kfOX
\kfOXItem 개미의 예측대로 며칠 뒤 큰비가 내렸다. \kfOX
\kfOXItem 성실하게 준비한 개미는 비가 와도 걱정이 없었다. \kfOX
\end{kfOXList}

\kfSecHeader{kfAreaLiterature}{문제}{}

\begin{kfQBlock}{01}{객관식}
\kfQStem{무더운 여름날, 베짱이는 무엇을 하고 있었나요?}
\begin{kfChoices}
\kfChoice 땀을 흘리며 먹이를 나르고 있었다.
\kfChoice 나무 그늘에서 바이올린을 켜고 있었다.
\kfChoice 집을 튼튼하게 수리하고 있었다.
\kfChoice 개미를 도와 식량을 모으고 있었다.
\end{kfChoices}
\end{kfQBlock}
\begin{kfQBlock}{02}{객관식}
\kfQStem{개미들이 더운 날에도 쉬지 않고 일한 이유는 무엇인가요?}
\begin{kfChoices}
\kfChoice 부지런함을 베짱이에게 자랑하고 싶어서
\kfChoice 너무 심심해서 마땅히 할 일이 없어서
\kfChoice 다가올 비와 겨울을 대비하기 위해서
\kfChoice 더운 날 땀 흘리며 운동을 하기 위해서
\end{kfChoices}
\end{kfQBlock}
\begin{kfQBlock}{03}{객관식}
\kfQStem{개미가 비가 올 것이라고 생각한 근거는 무엇인가요?}
\begin{kfChoices}
\kfChoice 베짱이가 슬픈 노래를 불렀기 때문에
\kfChoice 하늘을 보고 날씨를 예측했기 때문에
\kfChoice 부엌 달력에 비가 온다고 적혀 있어서
\kfChoice 이미 굵은 비가 내리고 있었기 때문에
\end{kfChoices}
\end{kfQBlock}
\begin{kfQBlock}{04}{객관식}
\kfQStem{이 이야기에서 얻을 수 있는 교훈으로 가장 알맞은 것은?}
\begin{kfChoices}
\kfChoice 노래를 잘 부르려면 연습을 많이 해야 한다.
\kfChoice 친구가 어려울 때는 모른 척해야 한다.
\kfChoice 미래를 예측하고 미리미리 준비하는 자세가 필요하다.
\kfChoice 여름에는 무조건 밖에서 놀아야 건강해진다.
\end{kfChoices}
\end{kfQBlock}
\begin{kfQBlock}{05}{서술형}
\kfQStem{베짱이에게 해주고 싶은 충고의 말을 한 문장으로 써 보세요. (힌트: '대비'라는 단어를 사용해 보세요.)}
\kfAnswerNote{답안}{4}
\end{kfQBlock}
\begin{kfQBlock}{06}{서술형}
\kfQStem{만약 베짱이가 일기 예보를 봐서 내일 비가 올 확률이 90\%라는 것을 알았다면 어떻게 행동했을까요? 상상해서 써 보세요.}
\kfAnswerNote{답안}{4}
\end{kfQBlock}
\begin{kfQBlock}{07}{서술형}
\kfQStem{개미의 말:}
\kfAnswerNote{답안}{4}
\end{kfQBlock}
\begin{kfQBlock}{08}{서술형}
\kfQStem{그 후 베짱이는:}
\kfAnswerNote{답안}{4}
\end{kfQBlock}
\begin{kfQBlock}{09}{서술형}
\kfQStem{내용:}
\kfAnswerNote{답안}{4}
\end{kfQBlock}


\kfStartNonlit{일기 예보는 어떻게 만들어질까요?}


\kfPassageNote{일기 예보는 어떻게 만들어질까요?}{%
　우리는 매일 아침 뉴스나 스마트폰으로 일기 예보를 확인합니다. "오후에는 비가 올 확률이 높으니 우산을 챙기세요."라는 말을 듣고 우리는 비에 대비합니다. 그렇다면 기상청에서는 어떻게 날씨를 미리 알 수 있을까요? 일기 예보는 과학적인 과정을 통해 만들어집니다.

　첫 번째 단계는 '관측'입니다. 현재의 날씨 상태를 정확하게 아는 것이 가장 중요하기 때문입니다. 땅에서는 기상 관측 장비로 기온, 습도, 바람 등을 재고, 바다에서는 기상 부표를 띄워 파도의 높이와 수온을 잽니다. 하늘에서는 '기상 레이더'가 구름 속에 비나 눈이 얼마나 들어있는지 살피고, 우주에 떠 있는 '기상 위성'은 지구 전체의 대기 흐름과 구름의 이동을 촬영하여 보내줍니다.

　두 번째 단계는 '분석'입니다. 관측을 통해 모은 엄청난 양의 정보를 '슈퍼컴퓨터'에 입력합니다. 슈퍼컴퓨터는 복잡한 수학 공식으로 지구의 공기 흐름을 계산하여 미래의 날씨 변화를 그려냅니다. 이것을 '수치 예보 모델'이라고 합니다.

　세 번째 단계는 '예보'입니다. 슈퍼컴퓨터가 계산한 결과를 바탕으로 날씨 전문가인 '예보관'들이 최종 결정을 내립니다. 예보관들은 컴퓨터의 계산 결과와 자신의 경험, 그리고 우리나라의 지형적인 특징 등을 종합하여 내일의 날씨를 판단합니다.

　이렇게 만들어진 일기 예보는 방송, 신문, 인터넷 등을 통해 우리에게 전달됩니다. 과학 기술의 발달과 전문가들의 노력 덕분에 우리는 변덕스러운 자연 현상에 미리 대처하고 안전하게 생활할 수 있습니다.
}


\kfActivityBg \par\kfMaybeNeedspace{32\baselineskip}\noindent{\kfTipMark\,\,}{\small\color{kfInk} 다음은 각 문단의 핵심 내용을 정리한 것입니다. 빈칸에 알맞은 말을 채워 넣으세요.}\par\nopagebreak[4]\smallskip\nopagebreak[4]
\begin{kfStructure}
\kfNode{0}{}{}
\kfNode{0}{}{}
\kfNode{0}{}{}
\kfNode{0}{}{}
\kfNode{0}{}{}
\end{kfStructure}

\kfActivityBg \par\kfMaybeNeedspace{24\baselineskip}\noindent{\kfTipMark\,\,}{\small\color{kfInk} 다음 관측 자료를 보고 내일의 날씨를 예측하여 방송 원고를 완성해 보세요.}\par\nopagebreak[4]\smallskip\nopagebreak[4]
\begin{kfBlankList}
\kfBlankItem [관측 자료] * 구름: 서쪽에서 먹구름이 빠르게 몰려오고 있음. * 바람: 바람이 점점 강하게 불고 있음. * 기온: 내일 아침 기온이 오늘보다 5도 정도 떨어질 것으로 보임.

[방송 원고] "내일 날씨를 알려드립니다. 서쪽에서 \kfFillBlank[12mm]{}이 몰려오고 있어 비가 내릴 것으로 보입니다. \kfFillBlank[12mm]{}도 강하게 불고 기온이 떨어져 날씨가 \kfFillBlank[12mm]{}지겠으니, 따뜻한 옷차림을 하시는 것이 좋겠습니다."
\end{kfBlankList}

\kfSecHeader{kfAreaNonlit}{문제}{}

\begin{kfQBlock}{01}{객관식}
\kfQStem{이 글의 중심 내용으로 가장 알맞은 것은?}
\begin{kfChoices}
\kfChoice 일기 예보에 쓰이는 슈퍼컴퓨터의 성능
\kfChoice 하늘에서 기상 위성이 하는 다양한 일
\kfChoice 일기 예보가 만들어지는 과학적 과정
\kfChoice 기상청에서 일하는 예보관들의 하루 생활
\end{kfChoices}
\end{kfQBlock}
\begin{kfQBlock}{02}{객관식}
\kfQStem{날씨를 예측하는 과정이 올바른 순서대로 나열된 것은?}
\begin{kfChoices}
\kfChoice 분석 → 관측 → 예보
\kfChoice 관측 → 분석 → 예보
\kfChoice 예보 → 관측 → 분석
\kfChoice 관측 → 예보 → 분석
\end{kfChoices}
\end{kfQBlock}
\begin{kfQBlock}{03}{객관식}
\kfQStem{다음 중 '관측' 단계에서 하는 일이 \uline{\underline{아닌}} 것은?}
\begin{kfChoices}
\kfChoice 기상 위성으로 구름의 이동을 촬영한다.
\kfChoice 슈퍼컴퓨터로 미래의 날씨를 계산한다.
\kfChoice 기상 레이더로 구름 속의 비를 살핀다.
\kfChoice 바다에 기상 부표를 띄워 수온을 잰다.
\end{kfChoices}
\end{kfQBlock}
\begin{kfQBlock}{04}{객관식}
\kfQStem{일기 예보를 만들 때 슈퍼컴퓨터를 사용하는 이유로 알맞은 것은?}
\begin{kfChoices}
\kfChoice 정보의 양이 너무 많고 계산이 복잡해서
\kfChoice 사람보다 글씨를 예쁘게 쓸 수 있어서
\kfChoice 인공위성을 조종하기 위해서
\kfChoice 전기 요금을 절약하기 위해서
\end{kfChoices}
\end{kfQBlock}
\begin{kfQBlock}{05}{객관식}
\kfQStem{"오후에 비가 올 확률이 80\%입니다"라는 일기 예보의 의미로 가장 적절한 것은?}
\begin{kfChoices}
\kfChoice 비가 80방울 떨어질 것이다.
\kfChoice 비가 올 가능성이 매우 높다.
\kfChoice 비가 오지 않을 가능성이 더 높다.
\kfChoice 비가 땅의 80\%만 적실 것이다.
\end{kfChoices}
\end{kfQBlock}


\kfStartWeekly{실력확인 영역 학습}


\noindent\textit{\small (제한시간: 30분 / 총 15문항)}\par\medskip
\par\noindent\textbf{[4-6] 다음 글을 읽고 물음에 답하세요.}\par\medskip
\kfPassageNote{지문 4}{%
햇볕이 쨍쨍 내리쬐는 무더운 여름날이었습니다. 베짱이는 나무 그늘 아래서 시원한 바람을 맞으며 바이올린을 켜고 노래를 불렀습니다. 그때, 개미들이 줄을 지어 땀을 뻘뻘 흘리며 지나갔습니다. 개미들은 자기 몸집보다 훨씬 큰 먹이를 짊어지고 부지런히 집으로 나르고 있었습니다. 베짱이가 개미를 보고 비웃으며 말했습니다. "이봐, 개미 친구들! 이렇게 화창한 날에 왜 그렇게 힘들게 일을 해? 나랑 같이 노래나 부르며 놀자." 맨 앞에 가던 개미가 잠시 멈춰 서서 땀을 닦으며 대답했습니다. "우리는 다가올 비와 겨울을 대비해서 식량을 모으는 중이야. 하늘을 보니 곧 큰비가 쏟아질 것 같거든."
}
\par\noindent\textbf{[7-9] 다음 글을 읽고 물음에 답하세요.}\par\medskip
\kfPassageNote{지문 7}{%
일기 예보는 과학적인 과정을 통해 만들어집니다. 첫 번째 단계는 '관측'입니다. 현재의 날씨 상태를 정확하게 아는 것이 중요하기 때문입니다. 땅, 바다, 하늘, 우주에서 다양한 장비로 날씨 정보를 수집합니다. 두 번째 단계는 '분석'입니다. 관측을 통해 모은 엄청난 양의 정보를 '( ㉠ )'에 입력합니다. ( ㉠ )은 복잡한 수학 공식으로 지구의 공기 흐름을 계산하여 미래의 날씨 변화를 그려냅니다. 세 번째 단계는 '예보'입니다. ( ㉠ )이 계산한 결과를 바탕으로 날씨 전문가인 '예보관'들이 최종 결정을 내립니다.
}
\begin{kfQBlock}{01}{객관식}
\kfQStem{다음 중 '대비'의 뜻으로 가장 알맞은 것은?}
\begin{kfChoices}
\kfChoice 지나간 일을 잊지 않고 깊이 후회함
\kfChoice 앞으로 일어날 일에 대해 미리 준비함
\kfChoice 어떤 일이 일어날 가능성을 계산함
\kfChoice 자연 현상을 기계로 자세히 관찰함
\end{kfChoices}
\end{kfQBlock}

\begin{kfQBlock}{02}{객관식}
\kfQStem{다음 빈칸에 들어갈 알맞은 말은?}
\begin{kfEvidence}"기상청에서 이번 주말에 많은 비가 내릴 것이라고 [　　　　　]했다."\end{kfEvidence}
\begin{kfChoices}
\kfChoice 예보
\kfChoice 관측
\kfChoice 습도
\kfChoice 확률
\end{kfChoices}
\end{kfQBlock}

\begin{kfQBlock}{03}{객관식}
\kfQStem{위 글에서 개미들이 무더운 날씨에도 열심히 일한 이유는 무엇인가요?}
\begin{kfChoices}
\kfChoice 노래를 못 불러서 심심했기 때문에
\kfChoice 베짱이에게 자랑하고 싶어서
\kfChoice 다가올 비와 겨울을 미리 준비하기 위해서
\kfChoice 집을 튼튼하게 짓기 위해서
\end{kfChoices}
\end{kfQBlock}

\begin{kfQBlock}{04}{객관식}
\kfQStem{베짱이가 개미들의 말을 듣고 보인 반응으로 알맞은 것은?}
\begin{kfChoices}
\kfChoice 개미들의 말이 맞다고 생각하여 함께 일을 도왔다.
\kfChoice 날씨가 좋으니 걱정하지 말라며 개미들을 비웃었다.
\kfChoice 비가 올 것을 걱정하여 우산을 준비했다.
\kfChoice 개미들에게 맛있는 먹이를 나누어 주었다.
\end{kfChoices}
\end{kfQBlock}

\begin{kfQBlock}{05}{객관식}
\kfQStem{다음 중 일기 예보가 만들어지는 순서로 올바른 것은?}
\begin{kfChoices}
\kfChoice 예보 → 관측 → 분석
\kfChoice 분석 → 예보 → 관측
\kfChoice 관측 → 분석 → 예보
\kfChoice 관측 → 예보 → 분석
\end{kfChoices}
\end{kfQBlock}

\begin{kfQBlock}{06}{객관식}
\kfQStem{일기 예보의 첫 번째 단계인 '관측'에서 하는 일이 \uline{\underline{아닌}} 것은?}
\begin{kfChoices}
\kfChoice 기상 위성으로 구름 이동 촬영하기
\kfChoice 기상 레이더로 비구름 관찰하기
\kfChoice 바다에 기상 부표 띄우기
\kfChoice 예보관이 최종 날씨 결정하기
\end{kfChoices}
\end{kfQBlock}

\begin{kfQBlock}{07}{객관식}
\kfQStem{다음 빈칸에 들어갈 알맞은 접속사는?}
\begin{kfEvidence}"기상청에서 태풍이 온다고 예보했습니다. [　　　　　] 어선들은 모두 항구로 피했습니다."\end{kfEvidence}
\begin{kfChoices}
\kfChoice 그러나
\kfChoice 또는
\kfChoice 왜냐하면
\kfChoice 따라서
\end{kfChoices}
\end{kfQBlock}

\begin{kfQBlock}{08}{객관식}
\kfQStem{'그러나(하지만)'와 같은 접속사를 사용하는 경우는 언제인가요?}
\begin{kfChoices}
\kfChoice 앞뒤 내용이 비슷하거나 이어질 때
\kfChoice 앞뒤 내용이 서로 반대될 때
\kfChoice 앞 내용이 원인이 되어 뒤에 결과가 올 때
\kfChoice 여럿 중에서 하나를 선택할 때
\end{kfChoices}
\end{kfQBlock}

\begin{kfQBlock}{09}{객관식}
\kfQStem{다음 중 밑줄 친 낱말이 바르게 쓰인 문장은?}
\begin{kfChoices}
\kfChoice 정답을 알아맞췄다.
\kfChoice 강풍에 문이 쾅 다쳤다.
\kfChoice 내일은 기온이 낮을 것입니다.
\kfChoice 닭이 알을 낮았다.
\end{kfChoices}
\end{kfQBlock}

\begin{kfQBlock}{10}{객관식}
\kfQStem{다음 문장의 괄호 안에 들어갈 알맞은 말은 무엇인가요?}
\begin{kfEvidence}"먹구름이 끼더니 [　　　　　] 비가 쏟아졌습니다." \\\relax
(지금 바로, 짧은 시간에라는 뜻)\end{kfEvidence}
\begin{kfChoices}
\kfChoice 금새
\kfChoice 금세
\kfChoice 틈새
\kfChoice 어느새
\end{kfChoices}
\end{kfQBlock}

\begin{kfQBlock}{11}{서술형}
\kfQStem{다음 설명에 해당하는 낱말을 쓰세요.}
\begin{kfEvidence}"얽혀 있는 것을 풀어서 자세히 그 내용을 밝힘"\end{kfEvidence}
\kfAnswerNote{답안}{4}
\end{kfQBlock}

\begin{kfQBlock}{12}{서술형}
\kfQStem{이 이야기에서 얻을 수 있는 교훈을 한 문장으로 쓰세요.}
\kfAnswerNote{답안}{4}
\end{kfQBlock}

\begin{kfQBlock}{13}{서술형}
\kfQStem{위 글의 빈칸 ㉠에 들어갈 알맞은 말을 쓰세요.}
\kfAnswerNote{답안}{4}
\end{kfQBlock}

\begin{kfQBlock}{14}{서술형}
\kfQStem{다음 문장에서 '누가(주어)'와 '어찌하다/어떠하다(서술어)'를 찾아 쓰세요.}
\begin{kfEvidence}"성실한 개미들은 추운 겨울을 대비하여 부지런히 먹이를 모았습니다."\end{kfEvidence}
\kfAnswerNote{답안}{4}
\end{kfQBlock}

\begin{kfQBlock}{15}{서술형}
\kfQStem{다음 괄호 안에 들어갈 알맞은 접속사를 쓰세요.}
\begin{kfEvidence}"여름에는 덥고 비가 많이 온다. [　　　　　] 겨울에는 춥고 건조하다."\end{kfEvidence}
\kfAnswerNote{답안}{4}
\end{kfQBlock}




\end{document}
